% =====================================================================
%  A Conditional Reduction of the 4D Pure Yang-Mills Mass Gap:
%  the reduction, the measurement, and the ledger
%
%  MERGED EDITION (v3) -- single self-contained paper.
%  Build:  pdflatex main.tex   (run twice)
% =====================================================================
\documentclass[11pt,a4paper]{article}

\usepackage[T1]{fontenc}
\usepackage[utf8]{inputenc}
\usepackage{lmodern}
\usepackage[margin=2.5cm,top=2.2cm,bottom=2.4cm,headsep=14pt]{geometry}
\usepackage{amsmath,amssymb,amsthm,mathtools,bm}
\usepackage{microtype}
\usepackage[table]{xcolor}
\usepackage{booktabs,array,tabularx,longtable}
\usepackage{graphicx}
\usepackage{tikz}
\usetikzlibrary{arrows.meta,positioning,calc,shapes.geometric,decorations.pathreplacing,patterns}
\usepackage{tcolorbox}
\tcbuselibrary{skins,breakable}
\usepackage{enumitem}
\usepackage{multicol}
\usepackage{caption}
\usepackage{titlesec}
\usepackage{fancyhdr}
\usepackage[colorlinks=true,linkcolor=blue!45!black,citecolor=blue!45!black,
            urlcolor=blue!45!black,bookmarksnumbered=true,bookmarksopen=true,
            bookmarksopenlevel=1,breaklinks=true,pdfborder={0 0 0}]{hyperref}

% ---------- colours ----------
\definecolor{cT}{HTML}{1D7A3A}
\definecolor{cC}{HTML}{1F5FA8}
\definecolor{cM}{HTML}{B45F06}
\definecolor{cN}{HTML}{A61B1B}
\definecolor{cE}{HTML}{6A3D9A}
\definecolor{cO}{HTML}{555555}

% ---------- label chips ----------
\newcommand{\labb}[2]{\fcolorbox{#1!35}{#1!8}{\textcolor{#1}{\textbf{\scriptsize #2}}}}
\newcommand{\LT}{\labb{cT}{T}}
\newcommand{\LC}{\labb{cC}{C}}
\newcommand{\LM}{\labb{cM}{M}}
\newcommand{\LN}{\labb{cN}{N}}
\newcommand{\LE}{\labb{cE}{E}}
\newcommand{\LO}{\labb{cO}{O}}

% ---------- claims ----------
\newcounter{claim}
\newcommand{\claimhead}[2]{%
  \refstepcounter{claim}%
  \textbf{#1~\theclaim}#2\hspace{0.35em}\ignorespaces}
\newtcolorbox{claimbox}[1]{enhanced, breakable=false,
  colback=#1!6, colframe=#1!75, boxrule=0pt, leftrule=2.6pt,
  arc=2pt, left=7pt, right=6pt, top=5pt, bottom=5pt,
  before skip=8pt, after skip=8pt, fontupper=\small,
  parbox=false}
\newenvironment{proofx}{\par\noindent\emph{Proof.}\ }{\hfill$\square$\par\medskip}
\newcommand{\src}[1]{\par\noindent\begingroup\sloppy\footnotesize\color{black!55}\textbf{Source.} #1\par\endgroup\vspace{4pt}}
\newtcolorbox{qbox}{enhanced, colback=black!3, colframe=black!35, boxrule=0pt,
  leftrule=2pt, arc=0pt, left=12pt, right=10pt, top=5pt, bottom=5pt,
  before skip=8pt, after skip=8pt, fontupper=\small, parbox=false}
\newtcolorbox{v3box}{enhanced, breakable, colback=cT!4, colframe=cT!55, boxrule=0pt,
  leftrule=2.6pt, arc=2pt, left=9pt, right=9pt, top=6pt, bottom=6pt,
  before skip=10pt, after skip=10pt, fontupper=\small, parbox=false}
\newtcolorbox{warnbox}{enhanced, breakable, colback=cN!4, colframe=cN!60, boxrule=0pt,
  leftrule=2.6pt, arc=2pt, left=9pt, right=9pt, top=6pt, bottom=6pt,
  before skip=10pt, after skip=10pt, fontupper=\small, parbox=false}

% ---------- helpers ----------
\newcommand{\code}[1]{\texttt{\footnotesize #1}}
\captionsetup{font=small,labelfont=bf,skip=5pt}
\captionsetup[table]{position=bottom,skip=4pt}
\captionsetup[longtable]{position=bottom,skip=4pt}
\setcounter{tocdepth}{3}
\setcounter{secnumdepth}{3}
\setlist[itemize]{leftmargin=1.4em,itemsep=1pt,topsep=3pt}
\setlist[enumerate]{leftmargin=1.6em,itemsep=1pt,topsep=3pt}
\allowdisplaybreaks
\setlength{\emergencystretch}{2em}

% ---------- section style ----------
\titleformat{\part}[display]
  {\centering\vspace*{0.4em}}
  {\normalsize\scshape\color{black!55}Part \thepart}
  {0.9em}
  {\LARGE\bfseries}
  [\vspace{0.75em}{\color{black!25}\rule{\textwidth}{0.5pt}}]
\titlespacing*{\part}{0pt}{0pt}{1.7em}
\titleformat{\section}{\large\bfseries}{\thesection}{0.6em}{}
\titleformat{\subsection}{\normalsize\bfseries}{\thesubsection}{0.55em}{}
\titleformat{\subsubsection}{\normalsize\bfseries\itshape}{}{0em}{}
\titlespacing*{\section}{0pt}{1.35em}{0.65em}
\titlespacing*{\subsection}{0pt}{1.0em}{0.45em}

% ---------- header ----------
\pagestyle{fancy}
\fancyhf{}
\renewcommand{\headrulewidth}{0.4pt}
\fancyhead[L]{\footnotesize\color{black!55}The reduction, the measurement, and the ledger (v3.1)}
\fancyhead[R]{\footnotesize\color{black!55}\thepage}
\setlength{\headheight}{13pt}

\hypersetup{pdftitle={A Conditional Reduction of the 4D Pure Yang-Mills Mass Gap: the reduction, the measurement, and the ledger},
            pdfauthor={Yuxuan Jiang},
            pdfsubject={Yang--Mills mass gap; conditional reduction; lattice measurement; negative survey},
            pdfkeywords={Yang-Mills, mass gap, Clay Millennium Problem, Kall\'en-Lehmann spectral
              measure, atom fraction, lattice gauge theory, strong coupling, cluster expansion,
              reflection positivity, no-go, Lean 4}}

\begin{document}

% =====================================================================
%  TITLE
% =====================================================================
\begin{titlepage}
\centering
\vspace*{2.2cm}

{\LARGE\bfseries A Conditional Reduction of the\\[0.35em]
4D Pure Yang--Mills Mass Gap\par}

\vspace{0.9em}
{\large\itshape The reduction, the measurement, and the ledger\par}

\vspace{0.5em}
{\normalsize Third Edition --- \textbf{merged}, superseding the $169$-page v2 and four sub-papers\par}

\vspace{1.1em}
{\color{black!30}\rule{0.58\textwidth}{0.6pt}}

\vspace{2.2cm}
{\large Yuxuan Jiang\par}
\vspace{0.4cm}
{\small Yang--Mills campaign\quad$\cdot$\quad Lean 4 formalisation $+$ lattice numerics\par}
\vspace{0.25cm}
{\footnotesize 24 November 2026\quad$\cdot$\quad v3.1\quad$\cdot$\quad Zenodo \texttt{10.5281/zenodo.20416516}\par}

\vspace{1.1em}
{\color{black!30}\rule{0.58\textwidth}{0.6pt}}

\vspace{1.7cm}
\begin{minipage}{0.86\textwidth}
\small
\textbf{Status of this document.} This is a conditional reduction and an honest audit of a
failed campaign. \textbf{The Clay mass gap is \emph{not} proved.} What is proved is a reduction,
a number of supporting theorems, and one measurement; what is not proved is stated as such, and
every route that was tried and did not work is listed with the exact reason it failed.

\vspace{0.5em}
\textbf{Grade symbols.} Every substantive claim carries one of:
\LT{}~\textbf{theorem} (proved; axioms at most \texttt{[propext, Classical.choice, Quot.sound]}) ;
\LC{}~\textbf{conditional} (a theorem given an explicit hypothesis, stated in the claim) ;
\LN{}~\textbf{no-go} (a theorem that a whole class of approaches cannot work) ;
\LM{}~\textbf{model} (verified in a solvable model or by exhaustive machine enumeration) ;
\LE{}~\textbf{measured} (numerical, not proved) ;
\LO{}~\textbf{open} (named, not closed).

\vspace{0.5em}
\textbf{Provenance rule.} Every bookkeeping number below carries the source file that produced it.
No claim is asserted without a file, and no number is asserted from memory.
\end{minipage}
\vfill
{\footnotesize\itshape
\textbf{2020 Mathematics Subject Classification.} 81T13 (primary); 81T25, 60K35, 47A10 (secondary).\par
\vspace{0.35em}
\textbf{Key words.} Yang--Mills mass gap; K\"all\'en--Lehmann spectral measure; atom fraction;
lattice gauge theory; cluster expansion; reflection positivity; negative results; Lean~4.\par}
\vspace{0.4cm}
\end{titlepage}

% =====================================================================
%  ABSTRACT
% =====================================================================
\section*{Abstract}
\addcontentsline{toc}{section}{Abstract}

We report the outcome of a sustained attempt to reduce, and then to close, the four-dimensional
pure $SU(N)$ Yang--Mills mass gap --- the question as posed in \cite{jaffe}. \textbf{We do not prove
the mass gap.} What we do establish is
four things, and one thing we measure.

\emph{First, a reduction.} The gap is written as a positivity statement at a single point at the
physical coupling, together with a monotonicity step in the volume and a tail-uniform floor in the
coupling. The rewriting is exact \emph{on the atom route} --- the class of arguments that certifies
the gap through the atom weight of one fixed operator --- and it is exact only there: outside that
class it is false, and two solvable models make it false (\S\ref{sec:survey-three}). The reduction is
itself the reason the problem does not move: every attempt to discharge one of its three legs turns
out to require the leg itself.

\emph{Second, what is proved.} A strong-coupling residue theorem with a corrected window
$\beta_A < 3.777\times10^{-3}$ (we previously reported $0.018394$, which is the $a=0$ activation root
and is \emph{illegal} as a criterion window --- an over-report by a factor $4.87$, corrected here);
a finite-support shadow of the isolation, verified as a Lean~4 proposition with no \texttt{sorry};
a uniqueness theorem for the target, conditional on the model specification and stated at grade
\LC{}; a no-go for the sixth species of proof, also at grade \LC{} rather than \LT{} because one
step does not survive a computability reading; and four machine-verified lemmas --- two
counterexamples to two natural implications, and two arithmetic certificates for the window.

\emph{Third, a measurement.} The atom fraction $w_1$ at $\beta = 2.5$ on $L_s = 6$ and $L_s = 8$ with
$L_t = 12$. For one operator $w_1 = 0.3427 \pm 0.0119$ ($28.8\sigma$); for another, on the same Markov
chain and the same $12000$ configurations, $w_1 = 0.0168 \pm 0.0142$. \textbf{The two differ by a
factor $20.2$.} This is not an error: it is the numerical realization of a one-line counterexample
($w_1 = \cos^2\theta$ at fixed spectrum) which shows that $w_1$ is a property of the pair $(H,O)$ and
not of $H$ alone. Any statement of the form ``$w_1 > 0$'' therefore has no content until the operator
$O$ is named.

\emph{Fourth, an exhaustive negative survey.} Twenty routes to uniformity in $(L, a)$ were attacked and
all fail. They fail with a uniform direction: nine give only \emph{upper} bounds on the quantity wanted
(compactness, reflection positivity, Berezin--Lieb, the probabilistic reading, the Krein--Rutman
residual, concavity, the RG modulus, two Ursell constants); nine are circular (the positivity
improvement, minimality, operator inequalities, definability, the strong--weak bridge, and the named
obstructions $\Delta_0$, $T4$, $H1$, $M(\beta)$); two have explicit counterexamples. Two produced
something usable: a block-Loewner inequality (which repairs a previously unproved step) and a sampler
bug whose correction invalidates four earlier numerical sign determinations.

\emph{Fifth, a ledger.} Thirty-eight definitional collisions, which are the project's own error
catalogue. We keep them in the paper because they are the mechanism by which the errors in this
paper were caught --- twice during the writing of this very manuscript.

\bigskip
\noindent\textbf{The verdict is negative and we state it plainly:} the campaign reduced the problem,
measured a quantity at the physical coupling, and mapped the wall. It did not break the wall.
The residual is not a gap in an argument; it is the absence of a method that produces lower bounds
rather than upper bounds, and every lower bound attempted here terminates in the target itself.

\clearpage
\tableofcontents
\clearpage

% =====================================================================
%  REVISION NOTE (v3.1)
% =====================================================================
\begin{v3box}
\textbf{Note on this revision (v3.1).} This edition changes bookkeeping, not grades --- with one
exception, item~R4 below, which is a grade that was over-stated and is now corrected downwards. No
claim is added. The corrections are listed here, in the spirit of the ledger of
\S\ref{sec:ledger}: each was found by a pass that did not write the sentence it corrects.

\begin{enumerate}[label=\textbf{R\arabic*.},leftmargin=2.8em,itemsep=2pt,topsep=3pt]
\item \textbf{Two over-report factors were mis-paired.} Against $3.777\times10^{-3}$ the factor for
      $0.018394$ is $4.87$, not $4.84$; $4.84$ belongs to the earlier $0.018269$
      (\S\ref{sec:est-window-fix}). The abstract, \S\ref{sec:verdict-new} and the two printed ratios
      now each name the number they divide.
\item \textbf{Leg (i) is a floor, not a sign.} v3 wrote it as $\forall(L,a)\,w_1>0$, which is
      strictly weaker than the target $\inf_{L,t}w_1>0$ of \S\ref{sec:survey-target}. The pointwise
      form is what the measurement reaches; the gap between the two is the content of
      \S\ref{sec:reduction-framework}. Both are now named separately, in the statement and in
      Table~\ref{tab:legs}.
\item \textbf{The reduction is marked as conditional on the atom route.} As printed in v3 the
      equivalence read as a statement about the gap \emph{tout court}, which \S\ref{sec:survey-three}
      contradicts: two solvable models have $w_1$ and $m$ pointing in opposite directions. The
      restriction is now inside the statement, and the abstract says so.
\item \textbf{$N09$-$1$ is graded \LC, not \LT.} Its own statement reads ``conditional on the $M6$
      specification and the $L\to\infty$ reading'', and the grading rule on the title page makes a
      theorem with printed hypotheses conditional. This is the one substantive grade change in the
      revision; \S\ref{sec:est-n09} records it where the statement is, and the places that quote the
      grade have been changed with it.
\item \textbf{Two of the four machine-verified lemmas are certificates, not counterexamples.}
      \S\ref{sec:est-counter} already said so in its body; the abstract and the subsection title now
      say so too.
\item \textbf{A notation clash between $m_1$ and $m_1^2$.} $\nu$ is carried on $m^2$ and
      $C(t)=\int e^{-mt}\,d\nu(m^2)$; the measurement section wrote both $d\nu(m)$ and
      $\nu(\{m_1^2\})$ in one display. A single notation paragraph in \S\ref{sec:reduction-equiv}
      now fixes the convention, and \eqref{eq:m-defs} follows it.
\item \textbf{Two numeral clashes in the prose.} The ``\#35/\#36'' of \S\ref{sec:verdict-ai} are
      corpus \emph{item} labels; they are rows~26/27 of Table~\ref{tab:ledger}, and are now written
      that way. The $0.25\%$ of ledger row~11 is $0.24\%$ (\S\ref{sec:survey-table}, row~7, where the
      same ratio is stated correctly).
\item \textbf{Section numbering.} The source restarted the section counter at Part~III, so v3 printed
      two \S3, two \S4, two \S5 and two \S6. Numbering now runs continuously through the five parts,
      and every cross-reference is a \verb|\ref| rather than a typed numeral.
\item \textbf{The bibliography was complete and unused.} v3 carried seventy-five entries and not one
      \verb|\cite| in the body. Citations have been added at the places where a result is
      \emph{inherited} rather than proved --- chiefly \S\ref{sec:reduction-inherited} and the route
      table of \S\ref{sec:survey-table} --- and the records were re-verified. Three records were
      wrong and are corrected in the list: the DOI of \cite{mrs} pointed at a fluid-dynamics paper,
      the year of \cite{gopfert} was 1982, and a volume of \cite{balaban} mixed a separate paper into
      the \emph{Large field} series. The phrase ``the log-Sobolev radius $1/48$'' has been replaced
      by what the source proves: $1/48$ is the \emph{edge of the window} in which the log-Sobolev
      and Poincar\'e inequalities hold.
\item \textbf{Presentation.} Every table now carries a number and a caption; the total-variation
      column of Table~\ref{tab:compact} states its normalisation ($L^1$, twice the usual); the
      bookmarks carry no raw macro; and the box colours are used for grades only --- the one box that
      had a grade chip without a graded claim (in \S\ref{sec:verdict-needed}) is now uncoloured.
\end{enumerate}
\end{v3box}
\clearpage

% =====================================================================
%  PARTS
% =====================================================================
% =====================================================================
%  PART I -- THE REDUCTION
%  (merged edition v3, section 1)
% =====================================================================
\clearpage
\part{The reduction}
\label{part:reduction}

This part does one thing: it writes the mass gap as an exact statement about a single number, and
then shows why that rewriting does not move the problem. The rewriting is exact; the obstruction is
that each of its three conjuncts, when one tries to prove it, requires itself. We state the
equivalence (\S\ref{sec:reduction-equiv}), identify the one framework in which the object can be
placed and show that everything that framework gives for free is a \emph{transport} of positivity
rather than a production of it (\S\ref{sec:reduction-framework}), take the three legs one at a time
and exhibit the circularity of each (\S\ref{sec:reduction-legs}), record the one sentence that
summarises the situation (\S\ref{sec:reduction-nomove}), and close with the qualification that
governs the whole part: $w_1$ is a property of the operator \emph{pair}, not of the Hamiltonian
(\S\ref{sec:reduction-operator}).

% ---------------------------------------------------------------------
\section{The equivalence}
\label{sec:reduction-equiv}

Fix the object once and for all. Let $O=\tfrac14\operatorname{tr}F^2$, let $\nu\ge 0$ be its
K\"all\'en--Lehmann spectral measure on $[0,\infty)$, and put
\begin{equation}
  M_0 \;:=\; \int d\nu \;=\; \operatorname{Var}(O),\qquad
  m_1^2 \;:=\; \inf\operatorname{supp}\nu,\qquad
  f_1^2 \;:=\; \nu(\{m_1^2\}),\qquad
  w_1 \;:=\; \frac{f_1^2}{M_0}\;\in\;[0,1].
  \label{eq:w1}
\end{equation}
So $w_1$ is the atom fraction at the bottom of the spectrum: the weight of the lightest excitation
in the source $O$, normalised by the total spectral weight. The bound $w_1\in[0,1]$ is an identity,
not an estimate: $0\le \nu(\{m_1^2\})\le\nu([0,\infty))=M_0$, for every positive measure, every
$(L,a)$, and every normalisation of $M_0$
(\code{research/v3/species/U01.md}, \S3). The K\"all\'en--Lehmann representation itself is not an
assumption of this paper but a theorem of the frame: $\nu\ge0$ and $f_1^2\ge0$ follow from
reflection positivity and the spectral theorem \cite{kl,os1973}.

\smallskip
\noindent\textbf{One notation, fixed once.} The variable carried by $\nu$ is $m^2$, so the atom of
the lightest state sits at the \emph{point} $m_1^2$ and is written $\nu(\{m_1^2\})$; the correlator
is then $C(t)=\int e^{-mt}\,d\nu(m^2)=\int e^{-\sqrt{s}\,t}\,d\nu(s)$ with $s=m^2$. Writing
$\nu(\{m_1^2\})$ and $d\nu(m)$ in the same display, as an earlier draft did, is a notational
collision of exactly the kind catalogued in \S\ref{sec:ledger}; the convention above is used
throughout.

\begin{claimbox}{cC}
\claimhead{Reduction}{.} With $w_1$ as in \eqref{eq:w1}, with $\beta=2.5$ the physical coupling of
this paper, and with the physical operator $O=\tfrac14\operatorname{tr}F^2$ held fixed,
\begin{align*}
  \textbf{Clay}\ \text{\emph{on the atom route}}\;\Longleftrightarrow\;&
  \underbrace{\inf_{L,t}\, w_1(L,2.5,t) \,>\, 0}_{\text{(i) the uniform floor at }\beta=2.5}
  \;+\;\underbrace{\exists\ \text{a monotone (or convergent) step in } L}_{\text{(ii) the volume step}}\\[2pt]
  &+\;\underbrace{\exists\,\Delta_0>0\ \ \forall \beta\ge\beta_0:\ m_{0^{++}}(\beta)\ge\Delta_0}_{\text{(iii) the }\beta\text{-tail-uniform floor}} .
\end{align*}
Three qualifications belong to the statement and are not commentary on it.
\emph{(a) Route.} ``On the atom route'' names the class of arguments that certifies the gap through
the atom weight of a fixed $O$; \S\ref{sec:survey-three} exhibits two solvable models in which $w_1$
and $m$ point in opposite directions, so the equivalence fails outside that class, and a proof
outside it cannot be obtained by refining (i)--(iii).
\emph{(b) Operator.} Leg (i) is stated at fixed $O$; \S\ref{sec:reduction-operator} shows that
dropping the qualification lets the quantity sweep all of $[0,1]$ at fixed spectrum.
\emph{(c) Floor, not sign.} Leg (i) asks for an infimum, not for pointwise positivity. The two differ
by exactly the compactness that \S\ref{sec:reduction-framework} shows one does not have, and it is
the pointwise form alone that the measurement of \S\ref{sec:measurement} supplies. The three
conjuncts are not three routes to one target: they are the three things a proof of the target must
supply, and the equivalence is a rewriting of the quantifiers, not a theorem about Yang--Mills.
\end{claimbox}
\src{Reduction and the three-leg split: \code{research/v3/species/U15.md}, \S1 and \S5 (the target
must be written as $w_1(\beta)\ge w_0(\beta)$ with $w_0(2.5)$ unknown); \code{U02.md}, \S5(c) (the
target is a full-family infimum, $\inf_{L,a,t}w_1\ge w_0$, not an endpoint value; at fixed $\beta$
the family is the $\inf_{L,t}$ written in the statement, since $a=a(\beta)$);
\code{UNIF3.md}, \S0 and \S4.1.}

The one conjunct that is not merely a restatement is (iii), and it is where the equivalence earns
its keep, because (iii) turns out to be the same statement as a continuity statement in the lattice
spacing.

\begin{claimbox}{cT}
\claimhead{Lemma}{.} (Tail identity.) Suppose $\beta\mapsto a(\beta)$ is strictly decreasing with
$a(\beta)\to 0$ (asymptotic freedom). Then for every $c_0>0$,
\[
  \exists a_0\ \forall a<a_0:\ m_{0^{++}}(a)\ge c_0\Lambda
  \quad\Longleftrightarrow\quad
  \exists\beta_0\ \forall\beta>\beta_0:\ m_{0^{++}}(\beta)\ge c_0\Lambda .
\]
Hence ``$a$-uniform'' and ``$\beta$-tail-uniform'' are not two requirements but one, written in two
units: $a<a_0\iff\beta>\beta(a_0)$, on the same quantity, in the same units.
\end{claimbox}
\begin{proofx}
Both sides are of the form ``$m\ge c_0\Lambda$ on a terminal set of a monotone parameter''. The map
$\beta\mapsto a(\beta)$ is a strictly monotone bijection onto its range, so it carries terminal sets
to terminal sets and preserves the inequality. No estimate on $m$ is used, and none is produced.
\end{proofx}
\src{\code{research/v3/species/UNIF3.md}, Lemma A (\S1.2); the arithmetic is the project's own, the
input ``$a\to 0\iff\beta\to\infty$'' is \code{research/uniqueness/REPORT.md}:43 as quoted there.}

The point of the lemma is negative and it is worth stating separately: it removes a whole class of
proposed second stages. There is no fifth direction called ``$a\to 0$'' that could be attacked after
the physical coupling has been settled, because a statement at a \emph{single} $\beta$ (including
$\beta=2.5$) constrains no sequence with $a\to 0$. What (iii) asks for is uniformity on a
\emph{tail}, and that is what (ii) and (i) also ask for, in their own variables.

\begin{table}[htbp]
\centering\small
\begin{tabularx}{\textwidth}{@{}l l l X@{}}
\toprule
Leg & Object & Status & Source \\
\midrule
(i)  & $\inf_{L,t}w_1(L,2.5,t)>0$ at fixed $O$ & \LO{} as a floor; \LE{} only pointwise & \code{RUN3.md}, \S3; \code{U11.md}, \S3.2 \\
(ii) & monotone/convergent step in $L$ & \LN{} for the present toolkit & \code{U02.md}, \S2; \code{U03.md}, \S2--\S3 \\
(iii) & $\beta$-tail-uniform floor $=\Delta_0$ & \LO{} (value); \LN{} (target eligibility) & \code{H1-FINAL.md}, \S0, \S3 \\
(iii)$'$ & ($a$-uniform form of (iii)) & \LT{} identical to (iii) & \code{UNIF3.md}, \S1.2 \\
\bottomrule
\end{tabularx}
\caption{The three legs, their objects, and their grades. Row (i) is the uniform floor; the pointwise
form, which is what a measurement can reach, is strictly weaker (\S\ref{sec:reduction-framework}).
The equivalence of rows (iii) and (iii)$'$ is the tail identity proved above.}
\label{tab:legs}
\end{table}

% ---------------------------------------------------------------------
\section{The unique framework}
\label{sec:reduction-framework}

Because $w_1$ is an atom mass, the object $\nu$ lives in a space where compactness costs nothing: a
variational upper bound (Ritz; the bottom of the spectrum is a minimum of a quadratic form
\cite{reed-simon}) gives $m_1^2\le\Lambda$, so $\nu/M_0$ is a probability measure on the
compact $K=[0,\Lambda]$, and $\mathcal{P}(K)$ is weak-$*$ compact and metrisable by
Banach--Alaoglu with the Prokhorov tightness hypothesis satisfied automatically. That is the whole
of the good news, and it is the reason the compactness class of arguments is out of play rather
than merely unpolished.

\begin{claimbox}{cT}
\claimhead{Theorem}{.} (Transport, not production.) Let $K$ be compact, $w_1:K\to[0,1]$ continuous,
and $w_1(x)>0$ for every $x\in K$. Then
\[
  \inf_{x\in K} w_1(x)\;=\;\min_{x\in K} w_1(x)\;=\;w_1(x_0)\;>\;0 ,
\]
and, conversely, $\inf_K w_1>0$ implies pointwise positivity on $K$. The two statements are
equivalent: compactness plus continuity adds nothing to pointwise positivity except the licence to
move the infimum to the point where it is attained.
\end{claimbox}
\begin{proofx}
A continuous function on a compact set attains its minimum (Weierstrass, one line; no Ascoli--Arzel\`a
is needed, since a single function is trivially equicontinuous). The converse is immediate. Hence
the two formulations are equivalent and the compactness hypothesis is not doing work.
\end{proofx}
\src{\code{research/v3/species/U01.md}, \S1.1--\S1.2 (``compactness does not produce positivity'' is
stated there as the whole content of the route); \code{U16.md}, \S3.2 for the choice of $K$ and the
freeness of tightness.}

So on a truncated window $K=\{L_1,\dots,L_K\}\times[a_{\min},a_0]$ with $a_{\min}>0$ the floor
$\inf_K w_1=w_1(L^*,a^*)>0$ is available, and it is an \emph{identity}: it carries a positivity that
was already assumed at every point of $K$ to a single point of $K$. It does not create positivity,
and it says nothing about $a_{\min}\downarrow 0$, because $a=0$ is not in $K$.

The second free item is the direction of semi-continuity, and its direction is the wrong one.

\begin{claimbox}{cT}
\claimhead{Proposition}{.} (Atoms are upper semi-continuous.) If $\nu_n\Rightarrow\nu$ weakly and
$x$ is fixed, then $\limsup_n \nu_n(\{x\})\le\nu(\{x\})$. Consequently, if the limit measure has no
atom at $m_1^2$, then $w_1^{(n)}\to0$: at the limit the atom can be destroyed, and upper
semi-continuity does not prevent it. A floor, by contrast, needs a \emph{lower}
semi-continuous bound.
\end{claimbox}
\begin{proofx}
By the Portmanteau theorem applied to the closed sets $\overline{B_\varepsilon(x)}$,
$\nu_n(\{x\})\le\nu_n(\overline{B_\varepsilon(x)})$, so
$\limsup_n\nu_n(\{x\})\le\nu(\overline{B_\varepsilon(x)})$; let $\varepsilon\downarrow0$.
\end{proofx}
\src{\code{research/v3/species/U01.md}, \S4.3, Proposition P-2; \code{U16.md}, \S0(c) and \S3.3.}

Weak-$*$ semi-continuity is therefore \emph{free in the direction ``finite volume $\to$ limit''} and
absent in the direction a floor requires, namely ``limit $\to$ nearby members''. The gap is not
technical: there is a compact family, with a limit point that carries an atom of mass $0.3$, whose
uniform floor is exactly $0$.

\begin{table}[htbp]
\centering\small
\begin{tabular}{@{}rrrrrr@{}}
\toprule
$h$ & mass & $\int s\,d\mu_h$ & $E=\text{bottom atom}$ & $\inf\operatorname{supp}$ & $\|\mu_h-\mu_0\|_{TV}$ \\
\midrule
$0.000$ & $1.00000$ & $2.40000$ & $0.300000$ & $1.00000$ & $0.000$ \\
$0.200$ & $1.00000$ & $2.40000$ & $0.000000$ & $0.80000$ & $0.600$ \\
$0.100$ & $1.00000$ & $2.40000$ & $0.000000$ & $0.90000$ & $0.600$ \\
$0.050$ & $1.00000$ & $2.40000$ & $0.000000$ & $0.95000$ & $0.600$ \\
$0.020$ & $1.00000$ & $2.40000$ & $0.000000$ & $0.98000$ & $0.600$ \\
$0.005$ & $1.00000$ & $2.40000$ & $0.000000$ & $0.99500$ & $0.600$ \\
$0.001$ & $1.00000$ & $2.40000$ & $0.000000$ & $0.99900$ & $0.600$ \\
\bottomrule
\end{tabular}
\caption{A compact family whose member at $h=0$ carries an atom of mass $0.3$ and whose uniform atom
floor is $0$. Total variation is quoted on the $L^1$ scale,
$\|\mu-\nu\|_{TV}=\int|d\mu-d\nu|$, so that the value $0.600$ is twice the usual normalisation
$\frac12\int|d\mu-d\nu|=0.300$; the constant offset is a convention, the vanishing of the floor is
not. Computed (\code{U16.md}, \S3.3).}
\label{tab:compact}
\end{table}

\noindent The family is $\mu_h=0.3\,\mathrm{Unif}[1-h,1+h]+0.7\,\delta_3$ on $\mathcal{P}([0,3])$,
with $\mu_0=0.3\,\delta_1+0.7\,\delta_3$. Every moment readout is constant to six digits
($\int s\,d\mu_h\equiv2.400000$), the family is compact, the limit point carries an atom of mass
$0.3$, and yet $\inf_{h\in[0,1/2]}E(\mu_h)=0$ while $\sup=0.3$ is attained only at $h=0$. Compact
family plus a limit atom does \emph{not} give a uniform atom floor (\code{U16.md}, \S3.3, machine
computed).

Finally, neither end of the parameter space is compact, and this is what removes the tame
theorems. At $a=0$ and at $L=\infty$ the family has \emph{no member}: the domain is open in $a$ and
the $L$'s are discrete with no interpolation (\code{U01.md}, \S2.3). The dimension grows like
$(L/a)^4\to\infty$ (\code{DF-CONTINUITY.md}:247, as quoted in \code{U15.md}), so no o-minimal or
other tameness hypothesis of finite dimension applies, and a uniformity theorem for a definable
family cannot be invoked even in principle. The honest formulation is that the obstruction is not
``not compact'' but ``not closed'': the family has no member at the point where its infimum is
sought.

\begin{claimbox}{cN}
\claimhead{No-go}{.} (The generator slot is empty.) The seven general tools for uniformity over an
infinite-dimensional family --- Banach--Steinhaus, Ascoli--Arzel\`a, Prokhorov, Baire category,
$\Gamma$-convergence, uniform large deviations, and induction on $L$ --- are all transports or
promotions: pointwise $\to$ uniform, local $\to$ global, weak $\to$ strong, $L\to L+1$. Not one of
them turns a non-uniform family into a uniform one. Banach--Steinhaus is the only one that imposes no
structure on the index set, and its hypothesis is literally ``the family is already uniform at every
point''; applied to $w_1$ it returns the trivial bound $1$.
\end{claimbox}
\src{\code{research/v3/species/U16.md}, \S1 (statements), \S2 (the table of what supplies
uniformity), \S4 (the $a\to0$ side is outside the range of all seven).}

% ---------------------------------------------------------------------
\section{Why each leg is circular}
\label{sec:reduction-legs}

\subsection{(i) Positivity needs the atom, and the atom is what is missing}

The leg asks for $w_1>0$, and $w_1$ is an atom weight. By the proposition above, the atom functional
is upper semi-continuous only, while a floor requires lower semi-continuity, and the failure is
realised: a thin atom $p_n\delta_{m_1^2}+(1-p_n)\delta_\Lambda$ with $p_n\to0$ has $w_1=p_n\to0$ with
$M_0$, $m_1^2$ and the gap all fixed (\code{Q6-H1-STATUS.md}:208--215), and there are gapped
atomless measures with $f_1^2=0$ and every moment, positivity and Hankel condition satisfied
(\code{R6-ATOM-ADVERSARIAL.md}:254--259). So compactness together with pointwise positivity does not
imply a uniform floor at the level of measures, and the disproof is by explicit witnesses rather
than by a failed search (\code{U01.md}, Proposition P-3).

The instruments do not repair this, and the reason is structural rather than quantitative. Every
instrument built in this campaign is a weak-$*$ continuous functional of $\nu$: the moments
$\int s^k\,d\nu$, the Ursell combination \code{$\phi^{U}$}, and the coefficient sequence
\code{$a_n$}. They are blind to atoms in a precise sense. For any weak-$*$
continuous $g$ one has $\sup_\nu|E(\nu)-g(\nu)|\ge f_1^2/2$, i.e.\ the error of the instrument is at
least the signal (\code{R3-ATOM-FUNCTIONAL.md}, Theorem R3-2, as recorded in \code{U15.md}, \S3);
with $M_0$ and the first moment fixed, $w_1$ ranges over an interval
(\code{WA-ALIGNMENT.md}:37--39, recorded in \code{U02.md}, \S6); with $(M_0,\chi)$ fixed, $f_1^2=1$
and $f_1^2=0.25$ are both
realisable (\code{LF-LASTFOUR.md}:221); and all of them --- $\operatorname{Var}$, Hessian, $E[H]$,
$\chi$, $f_1^2$ itself as a readout --- are translation-invariant normalised moments, hence
structurally blind to position (\code{LOOP-RUNBOOK.md}, \S69.1, the \code{tb01} verdict). The
variance families used to build them are the Efron--Stein ones \cite{es}, and the one lower-bound
mechanism of a different type, a transport--entropy criterion for a spectral gap \cite{vrs}, needs a
curvature input that the lattice measure at $\beta=2.5$ does not supply. Fourteen
named routes were checked against this target and none closes it: $0/14$, eleven no-go, two
conditional with the hypothesis equal to the conclusion, one numerical
(\code{research/v3/species/U15.md}, \S2). That is the circularity of leg (i): the positivity must be
produced by an instrument that can see an atom, and every instrument available is continuous in the
topology in which atoms are destroyed.

\subsection{(ii) Monotonicity is not there, and would not suffice}

The leg asks for a monotone step in the volume. There is no monotonicity theorem for $w_1(L)$, and no
$L$-sequence of measurements either; within the known families both directions occur. The
delocalised lowest mode gives $w_1=2/(L-1)$, i.e.\ $2,1,\tfrac23,\tfrac12,\tfrac25,\tfrac13,\tfrac27$
for $L=2,\dots,8$, strictly decreasing to $0$ (exact rational check, \code{U02.md}, \S2.2); a
localised light mode with finite correlation length gives a constant. The volume dependence of a
mass is in general a genuine $1/L$ expansion rather than a monotonicity \cite{luscher77}, which is
one more reason to expect no sign here. So $w_1(L)$ is not a property of the species, and no sign can
be extracted from the structure alone.

Even granting monotonicity, it does not give the leg, because a non-increasing sequence can decrease
to zero: $2/(L-1)\searrow0$, so ``non-increasing'' does not imply ``infimum positive''. And the one
positivity input available does not even supply monotonicity in the right family: reflection
positivity yields $m_{\mathrm{eff}}(t)\downarrow$, $\widehat{w}_1(T)\downarrow$ with
$\widehat{w}_1(T)\ge w_1$, and $\lambda_2^t\downarrow$ with $\lambda_2\le1$ --- all in the transfer
direction, all upper bounds on $w_1$, hence the opposite direction to a floor; in the spatial
direction reflection positivity yields existence, not monotonicity, and $w_1(L)$ is not a member of
its monotone family at all (\code{U03.md}, \S2--\S3). The leg is therefore not merely unproved: the
kind of object that would prove it does not exist in the toolkit.

\subsection{(iii) The tail floor is \texorpdfstring{$\Delta_0$}{Delta\_0}, and \texorpdfstring{$\Delta_0$}{Delta\_0} is dead}

The leg asks for a floor uniform on the coupling tail. That object is $\Delta_0$, and $\Delta_0$ has
already been judged: it is neither sufficient nor necessary for the gap, so it cannot serve as a
form of the target, and its value is open.

\begin{claimbox}{cN}
\claimhead{No-go}{.} ($\Delta_0$ is not a rewriting of the gap.) Two implications fail.
\emph{Insufficiency}: a thin atom has $\Delta_{\mathrm{dyn}}>0$ while $m=0$
(\code{Q6-H1-STATUS.md}:53--56, 110--114; machine-checked as \code{H1\_does\_not\_imply\_atom} with
no \code{sorry}). \emph{Unnecessity}: in two-dimensional Yang--Mills $\Delta_{\mathrm{dyn}}=0$ while
$m=3g^2/8>0$ (\code{N03-WFULL.md}, \S0; \code{H1-FINAL.md}, \S3.3). Hence
$\Delta_0>0\not\Rightarrow m>0$ and $m>0\not\Rightarrow\Delta_0>0$: the floor is neither a
sufficient nor a necessary form of the target.
\end{claimbox}
\src{\code{research/v3/species/H1-FINAL.md}, \S0 rows 4--5 and \S3.3; \code{N04-2D.md}:14,49 for
$G''=W_{\mathrm{geo}}+\Delta_{\mathrm{dyn}}$ with $W_{\mathrm{geo}}\ge2$ and
$\Delta_{\mathrm{dyn}}=O(e^{-2t})\to0$ in two dimensions.}

\smallskip
\noindent\textbf{Where the missing floor actually sits.} The operator whose spectral gap is wanted is
the Hessian of the action at a background, $G''=-D^{2}+2adF_B$, and its indefinite direction at a
constant chromomagnetic field is not an artefact of the expansion: the Savvidy--Nielsen--Olesen
background \cite{savvidy,no} carries a negative mode in the transverse-magnetic channel
\cite{leutwyler}, and the instability is physical rather than a gauge copy. This is why the coarse-mode
gap cannot be had from positivity of $G''$: $G''$ is not positive, so the measure that would supply
the floor sits partly on the negative side. What is wanted is therefore not a sharper estimate of
$G''$ but a statement about the spectrum of $\nu$ that does not route through $G''$ at all --- which
is leg (iii) again.

The solvable model $M6$ is where the insufficiency used to be argued, and that argument must now be
quoted with its quantifier attached: on $M6$ one has $\Delta_{\mathrm{dyn}}=2\omega$ on the
$a$-mode while $m_6=\min(\omega,2\pi^2/L^2)\to0$, but the $b$-mode gives $\Delta_{\mathrm{dyn}}=0$,
so $\inf_v\Delta_{\mathrm{dyn}}=0$ and $M6$ does \emph{not} satisfy the target under the
$\forall v$ reading of the quantifier (\code{H1-FINAL.md}, \S3.2, correcting
\code{N12-VERDICT.md}). The insufficiency witness is the thin atom; the unnecessity witness is two
dimensions.

What remains is a single open cell, and it is not a gap in an argument: it is the absence of a
$(L,a)$-uniform positive lower bound on $\Delta_{\mathrm{dyn}}$, equivalently a uniform spectral gap
of the coarse-mode operator $G''$ (\code{H1-FINAL.md}, \S0; \code{N11-FALLBACK.md}, \S6). That cell
is exactly the content of leg (iii), which is why the leg is circular: to prove it one must already
have it.

% ---------------------------------------------------------------------
\section{Why it does not move}
\label{sec:reduction-nomove}

\begin{qbox}
\textbf{The equivalence is itself the reason it does not move.} Each of the three legs, the moment
one sets out to prove it, turns out to require itself: (i) needs an atom seen by an instrument that
cannot see atoms, (ii) needs a monotone quantity that the positivity in hand does not supply and
that would not suffice if it did, and (iii) \emph{is} the missing cell, renamed.
\end{qbox}

This is worth separating from the ordinary experience of a hard problem, where a reduction shortens
the distance to the target. Here the reduction is exact and the distance is unchanged: the three
conjuncts are not three lemmas that would compose into a proof, but three descriptions of one
missing object, and the descriptions are equivalent precisely because they all name the same
uniformity. A reduction that is an equivalence of this kind cannot be a source of progress; it can
only be a source of bookkeeping, and the bookkeeping is what the rest of this paper performs.

% ---------------------------------------------------------------------
\section{The key qualification: \texorpdfstring{$w_1$}{w1} depends on the operator}
\label{sec:reduction-operator}

The last point is the one that changes how leg (i) must be written, and it is a two-line
counterexample rather than an empirical worry.

\begin{claimbox}{cT}
\claimhead{Theorem}{.} ($w_1$ is not a function of the spectrum.) Fix any Hamiltonian $H$ with
bottom $E_0$, first excitation $\psi_1$, second excitation $\psi_2$ and $\lambda_1=E_1-E_0>0$. For
$\theta\in[0,\pi/2]$ and $0<\varepsilon\ll1$ set
\[
  O_\theta \;=\; 1+\varepsilon\bigl(\cos\theta\,\psi_1+\sin\theta\,\psi_2\bigr),
\]
which is a legitimate non-negative source. Then $\operatorname{Var}(O_\theta)=\varepsilon^2$ and
$f_1^2=\varepsilon^2\cos^2\theta$, so
\[
  w_1(O_\theta)\;=\;\cos^2\theta\;\in\;[0,1],
\]
while $\operatorname{Spec}(H)$, $\lambda_1$, $L$ and $a$ are unchanged. Hence there is no function $F$
with $w_1=F(\operatorname{Spec}H)$; $w_1=f_1^2/M_0$ is a property of the pair $(H,O)$, not of $H$.
\end{claimbox}
\begin{proofx}
Insert $O_\theta$ into \eqref{eq:w1}. The variance is $\varepsilon^2$ by orthonormality of
$\Omega,\psi_1,\psi_2$; the bottom atom is the $\psi_1$ component squared, $\varepsilon^2\cos^2\theta$;
the ratio is $\cos^2\theta$. Nothing in the spectrum enters the computation.
\end{proofx}
\src{\code{U11.md}, \S3.2, equations (3.2)--(3.3), with the scope note there: this is a no-go on the
operator-pair class, not a claim that $w_1$ is spectrum-independent on Yang--Mills, where $(H,O)$ is a
single fixed pair (file: \code{research/v3/species/U11.md}).}

The counterexample is not a pathology of an unphysical source; it is realised in the measurement.
On one Markov chain, with the same $12000$ configurations, two operators give atom fractions that
differ by a factor $20.2$. That a mass estimate is an estimate \emph{for the operator that was
interpolated}, not for the theory, is the working premise of lattice spectroscopy
\cite{morningstar}; what the theorem above adds is that the dependence is total --- it can exhaust
the whole interval $[0,1]$ --- so no amount of statistics can remove it.

\begin{table}[htbp]
\centering\small
\begin{tabularx}{\textwidth}{@{}l r r r r r@{}}
\toprule
basis & $C(0)$ & $C(1)/C(0)$ & $m_{\mathrm{fit}}$ & $w_1$ (mass pinned) & $S/N(2)$ \\
\midrule
\code{raw}   & $19.888$ & $0.07372\pm0.00265$ & $1.707$ & $0.3427\pm0.0119$ & $5.0$ \\
\code{TAPE1} & $45.527$ & $0.00302\pm0.00297$ & $1.073$ & $0.0168\pm0.0142$ & $0.9$ \\
\code{TAPE2} & $39.428$ & $0.00150\pm0.00298$ & $0.762$ & $0.0107\pm0.0136$ & $1.6$ \\
\bottomrule
\end{tabularx}
\caption{One Markov chain, one set of configurations, three operators, all read at $m_1=1.546$
(Table~\ref{tab:m-bases}). Only the operator changes, and the atom fraction changes by a factor $20.2$.}
\label{tab:opdep-front}
\end{table}

\noindent Geometry $8^3\times12$, $\beta=2.5$, $\kappa=1.25$, $N=12000$ ($8$ chains $\times$ $1500$
configurations, $\mathrm{gap}=3$, $400$ thermalisation sweeps per chain), plaquette
$0.65234\pm0.00002$, errors from a $20$-block delete-one jackknife, $w_1$ with the mass pinned at
$m_1=1.546$. \code{TAPE} is a spatial APE smearing with the time direction left unsmeared
(\code{RUN3.md}, \S2--\S3). The two leading entries differ by the factor
$0.3427/0.0168=20.2$ reported in the abstract; the \code{raw} value reproduces an earlier
independent run at $0.02\sigma$.

Two consequences follow, and the second is a boundary of this paper.

\begin{v3box}
\textbf{(1) Leg (i) must be stated at fixed $O$.} ``$w_1>0$'' has no content until $O$ is named: at
fixed spectrum the value ranges over $[0,1]$. The correct reading of (i) is
\emph{for the physical operator $O=\tfrac14\operatorname{tr}F^2$, $w_1(2.5)>0$ at each $(L,a)$}, and
every statement in this paper is to be read that way.

\textbf{(2) The measured value is \LE{} and nothing more.} It is a lattice measurement at
$L_t=12$, $L_s=6$ and $L_s=8$, at one coupling, with $S/N(2)=5.0<10$: the $28.8\sigma$ quoted in the
abstract is significance from zero, not a demonstration of a floor, and the deficit in $S/N$ is
recorded rather than repaired (\code{RUN3.md}, \S4).
\end{v3box}

\begin{warnbox}
\textbf{Where the grades must not be guessed.} Three hooks in this part are grade-sensitive and are
flagged here rather than settled.

(a) The $M6$ row: the insufficiency claim is true per mode ($\Delta_{\mathrm{dyn}}=2\omega$ on the
$a$-mode) but false under the $\forall v$ quantifier ($\inf_v=0$), and the repository withdrew the
row on that ground (\code{H1-FINAL.md}:102,138). This part therefore uses the thin atom, not $M6$, as
the insufficiency witness; any sentence restoring $M6$ to that role would need a grade that is not
available in the sources.

(b) The factor $20.2$ compares two \emph{operators}, not two theories: it is evidence for the
operator dependence, and it is not evidence for or against any floor.

(c) Leg (i) at fixed $O$ and the mass gap as usually stated differ by an $O$-independence
requirement, and no such statement exists in this paper. That is a genuine boundary of the
reduction, not a technicality: the reduction is exact at fixed $O$, and the passage to a statement
about $H$ alone is exactly what the counterexample forbids.
\end{warnbox}

% ---------------------------------------------------------------------
\section{Where the reduction sits in the literature}
\label{sec:reduction-inherited}

Almost nothing in this part is new except the reduction itself, and a reader is entitled to know which
sentences carry an obligation and which are borrowed. The following inventory is organised by what is
borrowed, not by chronology; every item is used somewhere in the paper, and none of them is used at a
grade higher than the source supports.

\noindent\textbf{(a) The frame.} The problem is the one posed in \cite{jaffe}: a quantum field theory
satisfying the Wightman axioms \cite{sw}, with a mass gap above the vacuum. The spectral data
$\nu$, $m_1$, $f_1^2$ of \eqref{eq:w1} are the K\"all\'en--Lehmann data of that frame \cite{kl}, and
the positivity that makes $\nu$ a measure at finite volume is reflection positivity
\cite{os1973}, in the lattice form of Osterwalder--Seiler \cite{osterwalder-seiler}.

\noindent\textbf{(b) The strong-coupling mass gap is not ours.} That lattice gauge theory has a mass
gap at large $\beta$ is a theorem of the 1970s and 1980s: Seiler's lecture notes \cite{seiler}, the
Brydges--Giffen--Durhuus--Fr\"ohlich polymer expansion \cite{bgdf}, and, for the lower-dimensional
models, Glimm--Jaffe--Spencer \cite{gjs} and Feldman--Osterwalder \cite{fo}. The criterion that the
five rings of \S\ref{sec:est-i1} are run against is Koteck\'y--Preiss \cite{kp}; the edge $1/48$ of
the strong-coupling window is the threshold of Shen--Zhu--Zhu's log-Sobolev and Poincar\'e
inequalities \cite{szz}. None of
these is reproved here. What \S\ref{sec:est-i1} adds is the residue statement $I1$ and, with it, the
removal of a premise that the earlier literature needed --- that the $0^{++}$ channel be strictly the
lightest.

\noindent\textbf{(c) The renormalisation-group route.} The programme that would prove the gap at all
couplings by a multiscale expansion is Balaban's \cite{balaban,balaban-conv}, and its completion with an infrared
cutoff is Magnen--Rivasseau--S\'enor \cite{mrs}. Neither is used in this paper, and the reason is
recorded rather than hidden: the estimates are not uniform in the ways \S\ref{sec:survey} needs, and
the constants in them are not the constants this campaign can control.

\noindent\textbf{(d) Analytic tools.} Log-concavity enters through Pr\'ekopa \cite{prekopa} and
Brascamp--Lieb \cite{bl}; the spectral statements through Kato \cite{kato} and Reed--Simon
\cite{reed-simon}; level crossings through von Neumann--Wigner \cite{vnw}; the operator-comparison
bounds of row 12 of \S\ref{sec:survey-table} through Temple \cite{temple}; the stochastic and
semigroup side through Feller \cite{feller}, Jerrum--Sinclair \cite{js}, Sinclair \cite{sinclair},
Bauerschmidt--Brydges--Slade \cite{bbs} and the lecture notes of \cite{chatterjee}.

\noindent\textbf{(e) Other routes to a mass scale, and why none of them is the missing machine.}
Lattice formulations and their numerical study are Wilson's \cite{wilson} and Creutz's \cite{creutz},
with the step-scaling and flow technologies of \cite{luscher-stepscaling,luscher-flow} and the
spectroscopy of \cite{morningstar}; the large-$N$ and loop-equation route is
\cite{thooft-planar,eguchi-kawai,makeenko-migdal}; the confining-string route is
\cite{polchinski-strominger,luscher-weisz}; the holographic route is
\cite{maldacena,gkp,witten,witten-thermal,polchinski-strassler,sakai-sugimoto}; the conformal-bootstrap
route is \cite{elshowk,prv}; the resurgent route is \cite{abs}; the $\beta$-function and
$a$-theorem complex is \cite{caswell,vrvl,banks-zaks,polchinski,komargodski-schwimmer,dksst,zamolodchikov,osborn},
with the conformal-collider bounds of \cite{hofman-maldacena,hofman-li} and the near-conformal lattice
studies of \cite{aoki}; and the infrared scenarios that put the mass scale in the gauge sector are
Gribov's \cite{gribov}, its topological formulation by Singer \cite{singer} and Zwanziger
\cite{zwanziger}, the Dyson--Schwinger truncations of \cite{vsmekal,fischer-alkofer} reviewed in
\cite{huber}, and the dynamically generated gluon mass of \cite{cornwall}. Confinement itself has
been proved in a restricted model, three-dimensional $U(1)$ \cite{gopfert}.

Each entry in (e) supplies an upper bound, a scenario, or a numerical value at a fixed coupling. None
of them supplies the object \S\ref{sec:verdict-needed} names, and the uniform direction of
\S\ref{sec:survey} is the observation that this is not an accident of any one route.

\label{part:reduction-end}

% =====================================================================
%  sec02-established.tex
%  Part II: what is established.
%  Sources read for this section:
%    research/v3/lean/{M02-CBAR, L04PropC, L05Counterexamples, M03-COUNTER}.lean
%    research/v3/species/{C06-WINDOW, ADJ2, ADJ4, C07-AUDIT, F07-ADJCBAR,
%      D05-FKG, N09-RESTATE, L02-CBARKATO, L06-INTERVAL, Y01-KL,
%      Z06-HOLEA, Z07-HOLEB, Z08-HOLEC, Z09-HOLED, Z10-HOLEE, Z12-FIXC}.md
%    docs/paper-v2/ch20-unconditional-i1.tex (five rings, five holes, window)
%    docs/paper-v7/sg1/main.tex and docs/paper-v7/w1/main.tex
%    docs/position-current.html:373--417 (round-178 status block)
%  House rules honoured here: every number carries a source file; grades are
%  never written above what the sources support; no completion stamp is
%  printed anywhere; Clay = 0%.
% =====================================================================
\clearpage
\begingroup
\renewcommand{\code}[1]{\texttt{\footnotesize\detokenize{#1}}}

\part{What is established}


\section{What is established: the residue, the isolation, the uniqueness, the no-go, and four machine-verified lemmas}
\label{sec:established}

\noindent This section is the register of statements that survived audit. It is deliberately narrow:
every claim below carries a grade chip and the file that produced it, and no number is quoted from
memory. Of the six subsections, three are theorems in the strict sense (\LT: $I1$, $I2$, and the
block-Loewner order), two are theorems under hypotheses printed with them (\LC: $N09$-$1$ and
$SG1$-$1$), one is model-level (\LM: the four machine-verified lemmas), and two carry corrections that
reduce what an earlier edition of this work asserted --- the window of \S\ref{sec:est-window-fix} and
the grade of \S\ref{sec:est-sg1-grade}.

\smallskip
\noindent\textbf{Two grading rules, used throughout.} (1) A claim is graded by the \emph{weakest}
link of the chain that certifies it, not by the strength of its conclusion. (2) A \emph{computed}
number never upgrades a \emph{proved} structure, and a proved structure never upgrades a computed
number: in \S\ref{sec:est-i1} the two are physically separated inside one theorem, and that
separation is the point of the subsection. \src{\code{position-current.html:373--380}.}

\noindent\textbf{Scope, stated once and not repeated.} Nothing in this section crosses to
$\beta\approx2.5$, nothing takes $a\to0$, and no statement below mentions the continuum mass gap.
The residue theorem is a strong-coupling statement at fixed spacing; its window is in the coupling,
and the spacing has a positive lower bound inside it (\code{docs/paper-v2/ch20-unconditional-i1.tex}
\S D.3). \textbf{Clay remains $0\%$.}

% =====================================================================
\subsection{The strong-coupling residue theorem \texorpdfstring{$I1$}{I1}: \texorpdfstring{\LT}{[T]} on the structure, \texorpdfstring{\LC}{[C]} on the value}
\label{sec:est-i1}

\begin{claimbox}{cT}
\claimhead{Theorem}{(strong-coupling residue $I1$).} Let
$\mathcal C_L=T^3_L\times_a\mathbb Z$ be the cylinder (spatial torus of side $L$, discrete time of
spacing $a$), let $O=\tfrac14\operatorname{tr}F^2$, and let $m_1$ be the lightest energy carrying
nonzero $O$-overlap. Fix $N$ and work in the SZZ coupling $\beta_A=2/(g^2N)$ (Wilson:
$\beta_W=N^2\beta_A$). For every finite $L$, \emph{uniformly in $L$},
\begin{equation}
\beta_A\ <\ \beta_\star\ :=\ W(a^\ast)\ =\ 3.777\times10^{-3}\ =\ 0.1813\times\tfrac1{48},
\qquad\text{i.e.}\qquad
\beta_W\ <\ 0.015108\ =\ 0.725\times\tfrac1{48},
\label{eq:est-window}
\end{equation}
the lightest $0^{++}$ atom of $O$ obeys
\begin{equation}
f^2_{1,L}\ :=\ \nu_L(\{E_{1,L}\})\ =\ \lim_{t\to\infty}e^{E_{1,L}t}C_L(t)\ \ge\ r_0\ >\ 0 .
\label{eq:est-f1}
\end{equation}
The order of limits is: $L$ finite and fixed, then $t\to\infty$; $L$ is a uniformity axis and
$L\to\infty$ is never taken. \emph{Structure: \LT. The window \eqref{eq:est-window}: \LC, computed
and not proved} (\S\ref{sec:est-window-fix}).
\end{claimbox}

\noindent\textbf{The five rings, all \LT.} The chain to \eqref{eq:est-f1} has five links. Each is one
line here because the work is in the file; naming them is what makes the grade auditable.

\begin{enumerate}[label=\textbf{R\arabic*.},leftmargin=2.7em]
\item \textbf{Four walls} (\code{P3-FOURWALL-DERIVATION.md} \S1.1--1.5). Each of the source
      plaquette's four boundary links carries an open fundamental index pair and $\int dU\,U_{ij}=0$
      forbids covering it alone; two unit plaquettes share at most one link, so there are at least
      four independent strands; only temporal plaquettes advance time, so each strand crosses $t/a$
      slices. Hence $|P|\ge4t/a$, with equality exactly for the straight four-walled tube, whose
      weight is $a_0u^{4t/a}=a_0e^{-m_1t}$ with $m_1=4\ln(1/u)/a$. \LT.
\item \textbf{KP in the full representation} (\code{P3-FULLREP-KP.md} \S2.1--2.3, \S4.1). The
      criterion is Koteck\'y--Preiss \cite{kp}. Haar is a
      probability measure, so $|\chi_r(U)|\le\dim r$; the multiplicity cap
      $c_r\le\beta^{|r|}/(|r|!\dim r)$ cancels exactly the $\dim r$ that the contraction would cost;
      the factorial beats the sub-exponential rep-label entropy; hence $\sum_{P\ni x}|w(P)|e^{a|P|}
      \le a$ whenever $uK<1$, with $K\le6.1\times10^2N$ and plaquette adjacency
      $\Delta=4(6-1)=20$. \LT.
\item \textbf{Resummation of the decorations} (\code{X01-L1C.md} \S1--2, \code{Y04-AUDITX01.md}
      \S1--3; activities corrected in \code{Z08-HOLEC.md}, \code{Z12-FIXC.md}). Clusters are
      grouped by attachment span $\alpha$, not by support; the tube is open, so placements are
      counted by $(T-\alpha+1)_+$; the shift $\alpha\mapsto\alpha-1-T$ moves weight and exponent
      together, so the tail is $|\varepsilon_T|\le B\bar u^T$ with no $T$-prefactor; the
      interface-corrected activity has envelope $\hat C=4$. \LT.
\item \textbf{Spectral bridge} (\code{P3-CYLINDER-RESIDUE.md} \S1.2--1.4). At finite $L$ the
      transfer matrix is Hilbert--Schmidt on a compact configuration space, hence compact, hence
      purely discrete: every spectral point is an atom, and $f_1^2=\lim_te^{E_{1,L}t}C_L(t)$ uses
      only $\delta_L=E_{2,L}-E_{1,L}>0$ at that $L$ \cite{reed-simon}. \LT.
\item \textbf{$[KL]$} (\code{P3-CYLINDER-RESIDUE.md}:67--69, \code{Y01-KL.md} \S4.1).
      $\nu_L\ge0$ is Osterwalder--Seiler reflection positivity \cite{osterwalder-seiler,os1973} and
      $f_{n,L}^2\ge0$ is the spectral theorem, so the K\"all\'en--Lehmann representation \cite{kl} on
      the cylinder is a theorem with no Yang--Mills dynamics in it, and it transfers verbatim to the
      $2^{++}$ channel. \LT.
\end{enumerate}

\noindent\textbf{The five holes A--E, all closed.} Each hole is a specific defect that an earlier
draft absorbed by relabelling; each is now closed in its own file.

\begin{table}[htbp]
\centering
\footnotesize
\setlength{\tabcolsep}{4pt}
\renewcommand{\arraystretch}{1.2}
\begin{tabularx}{\textwidth}{@{}l>{\raggedright\arraybackslash}p{3.1cm}%
  >{\raggedright\arraybackslash}X>{\raggedright\arraybackslash}p{3.3cm}@{}}
\toprule
\textbf{Hole} & \textbf{What it was} & \textbf{How it is closed} & \textbf{Source} \\
\midrule
A & $u=c_f$ vs.\ $u=c_f/N$: a factor $N^{4T}$ in the tube weight
  & $u:=c_f/N$, one $1/N$ per tile, from Gupta (14.18)/(14.21); and $A^{\mathrm{dressed}}$
    \emph{is} $f_1^2$, so no second quantity ever existed
  & \code{Z06-HOLEA.md} \S1--5 \LT \\
B & whether $\delta\ge\kappa/3$ lands inside $I1$
  & out: neither $I1$ nor $D_\infty$ contains $\delta$; the reachable rate is
    $\delta=\ln(1/\bar u)/a=\kappa/4-\ln(\hat Ce\Delta\bar C)/a$, strictly below $\kappa/4$
  & \code{Z07-HOLEB.md} \S2--4 \LT \\
C & re-coupling is not block-multiplicative: $\mathrm{Hom}(1,f\bar ff\bar f)=2\ne1^2$; caps have
    $w(\gamma)=0$ while $w(B_T\cup\gamma)\ne0$
  & $\hat C=2$ exactly at each factorization point, envelope $\hat C=4$ per plaquette; the
    interface re-coupling is absorbed into the activity, never into $\varepsilon_T$
  & \code{Z08-HOLEC.md} \S2--4, \code{Z12-FIXC.md} \S1--3 \LT \\
D & summing $O$ over the slice makes $f_1^2\propto L^3$, so a bare $f_1^2$ cannot meet an
    $L$-independent $r_0$
  & $r_0$ carries the volume factor (equivalently $f_1^2$ is read normalized); $a_0$'s
    $L$-independence is already proved by a positive geometric sum with wrap corrections $e^{-cL}$
  & \code{Z09-HOLED.md} \S3--5 \LT \\
E & $r_0=\tfrac14(2N)^{-2}$ vs.\ $\tfrac12c_1u^p$: two $\beta$-scalings for one constant
  & $p$ is the operator-coefficient exponent and $p=0$; $u$ lives only in the correction factor,
    $r_0=\tfrac12a_0(1-O(u))$, $\beta$-independent
  & \code{Z10-HOLEE.md} \S3--5 \LT \\
\bottomrule
\end{tabularx}
\caption{The five holes A--E of the residue-and-isolation chain, each with the defect that it was and
the argument that closes it. Every row is \LT{} in its own file; the sources are named in the last
column.}
\label{tab:holes}
\end{table}

\noindent\textbf{The two accounting constants.} \emph{$\bar C\equiv N=2$, as an identity and not as
an upper bound.} The Ursell constant is $\bar C=c_f/u=N$, one $1/N$ per link from the integration,
with $|w|^{1/n}/u=c_f/u=N$ on cubes and four-walled tubes (measured $2.0000$ at $n=6,4,8,12$); the
$\mathrm{SU}(2)$ cross-check at $\beta=1/12$ is $\langle\tfrac1N\operatorname{tr}U_p\rangle
=0.0208273$ against $c_f=0.0208333$, a $0.03\%$ gap.
\src{\code{F07-ADJCBAR.md} \S3; \code{M02-CBAR.lean}, \code{cbar_eq_N}, \code{cbar_eq_N_two}.}

\smallskip
\noindent\emph{$\hat C=4$, and attained rather than bounded.} The KP criterion sums over polymers
$\gamma$, so what enters is the per-polymer rate $\sup_\gamma(\prod_\ell R(\ell))^{1/|\gamma|}$; the
handshake identity $E_2/F=2$ on closed clusters gives $2^{E_2/F}=2^2=4$, and a witness attains it:
all $324/324$ cubes of $T^4_3$ give geometric mean $4.000000$, as do $4000$ random closed
complexes, with $0$ violations. The "$2$" is the value at each factorization point, not the value on
a polymer, so it can never replace the $4$. \LT{}$+$\LE.
\src{\code{ADJ2.md}:13; \code{C06-WINDOW.md} \S2.2--2.3.}

\smallskip
\noindent\emph{The two constants meet the threshold.} $\hat C\bar C=8$ against
$2N^2\cdot48/e\Delta=7.0632853\ldots$, and the Lean file contains the integer certificate for
$8>7.0632853$. \src{\code{M02-CBAR.lean}: \code{chat_times_cbar}, \code{thresh},
\code{thresh_iff}, \code{thresh_eq_eight_div_x}.}

\smallskip
\noindent\emph{Penrose, approved with its own caveat.} The tree bound $|\varphi^U(\Gamma)|\le
\tau(G_\Gamma)$ is \LT{} with a four-line deletion--contraction proof plus $875$ graphs and $0$
violations; the \emph{citation} to Penrose (1967) and to the Penrose identity in
Fern\'andez--Procacci stays \LC{} because the original page was never opened
(\code{C07-AUDIT.md}:17, :24, :124--126, :135).

\smallskip
\noindent\emph{Lean.} The arithmetic core of the two constants is $25$ theorems in
\code{M02-CBAR.lean} with \code{sorry} $=0$ and twelve axiom-audit lines returning
$[\code{propext},\code{Classical.choice},\code{Quot.sound}]$; the last run exits $0$ with no error
and no warning. \src{\code{M02-CBAR.lean} (theorem count and axiom audit at its end);
\code{LOOP-RUNBOOK.md}:9512--9528.}

% ---------------------------------------------------------------------
\subsubsection{The window correction: the value drops to \texorpdfstring{\LC}{[C]}, the structure does not drop}
\label{sec:est-window-fix}

\begin{warnbox}
\textbf{What was over-reported, by how much, and why the old number is illegal.} The window printed
in the earlier edition of this chapter was $\beta<0.018394$. That number is
$1/(e\Delta)$, the \emph{activation root at $a=0$} of the KP criterion --- and $a=0$ sets the
right-hand side of the criterion to zero, so the criterion is identically false there and the root
is a convergence radius, not a criterion window. \textbf{It is therefore not a legal window.} With
the grading parameter optimized,
\begin{equation}
\beta_A\ <\ W(a^\ast)\ =\ 3.777\times10^{-3}\ =\ 0.1813\times\tfrac1{48},
\qquad\Longleftrightarrow\qquad
\beta_W\ <\ 0.015108\ =\ 0.725\times\tfrac1{48},
\label{eq:est-window2}
\end{equation}
so the earlier report over-states the window by
\[
\frac{0.018269}{0.003777}\ =\ 4.84\times ,
\qquad\text{and this chapter's own}\quad
\frac{0.018394}{0.003777}\ =\ 4.87\times .
\]
\src{\code{ADJ4.md}:16 (the ratio $1:2.7508:7.7238$ and $\Phi$ table), :19 and :125 (the two
over-report factors); \code{LOOP-RUNBOOK.md}:17133--17194 (\S131, the same two factors);
\code{C06-WINDOW.md} \S2.}
\end{warnbox}

\noindent\textbf{The optimizer, printed so that it can be checked.} With $\lambda=e\Delta=54.3656$
and $\Phi$ the single-plaquette sum, the criterion $S_x(a)\le a$ is equivalent to
$\lambda\Phi(a,\beta)\le a/(1+a)$, hence $W(a)=a/[(1+a)\lambda\Phi(a)]$. The measured ratios are
$\Phi(0):\Phi(1):\Phi(2)=1:2.7508:7.7238$ at $\beta=0.018394$ ($2.75512$ at $\beta=1/48$), and the
numerical maximum sits at $a^\ast=0.6150$. The closed-form point $a^\ast=(\sqrt5-1)/2=0.618$ is an
approximation to that maximum, with the same value to four digits, and $a=1$ gives
$3.372\times10^{-3}$ (the $3.357\times10^{-3}$ of the earlier audit, $0.44\%$ lower).
\src{\code{ADJ4.md}:16, :19; \code{CB4.md}:53--57 as quoted there; \code{LOOP-RUNBOOK.md}:17149.}

\smallskip
\noindent\textbf{And the natural coupling is outside, by a wide margin.} At
$\beta_W=1/12=0.08333$ one has $u=\beta_W/(2N^2)=0.0104$ and $\bar u=e\Delta N u=1.132617>1$, so the
criterion fails; against the corrected window the same coupling is out by
$0.08333/0.015108=5.5\times$. \src{\code{M02-CBAR.lean} (\code{cf} at $\beta=1/12$ gives
$c_f=1/48$); \code{BETA-CONVENTION.md}:187; \code{F07-ADJCBAR.md} \S4.}

\begin{qbox}
\textbf{Why this is a correction of a value and not a retraction of a theorem.} The window is one
slot of the statement \eqref{eq:est-window} and it is not a premise of any ring of the chain: the
five rings \LT{} deliver $f^2_{1,L}\ge r_0>0$ for every $\beta$ at which the KP criterion holds, and
``at which it holds'' is exactly what the window names. Re-optimizing the grading parameter changes
\emph{where} the criterion holds, not \emph{what} it delivers. Hence the grade splits: the value of
the window is \LC{} (computed: a numerical maximum of an explicit function), the residue statement
stays \LT. \textbf{There is no conflict between the correction and the \LT{} structure, because the
two occupy different slots of the same sentence}; writing the whole theorem \LC{} would under-state
the chain, and writing the window \LT{} would over-state a numerical maximum.
\src{\code{ADJ4.md}:19 (the verdict separates the number from the reason); \code{position-current.html}:375.}
\end{qbox}

% =====================================================================
\subsection{The isolation \texorpdfstring{$I2$}{I2} and Proposition C}
\label{sec:est-i2}

\begin{claimbox}{cT}
\claimhead{Theorem}{(isolation $I2$).} In the setting of \S\ref{sec:est-i1}, with
$\Lambda_\star:=\inf\operatorname{supp}\,(\nu-f_1^2\delta_{m_1^2})$ the bottom of the spectrum above
the leading atom, for every finite $L$ and uniformly in $L$,
\begin{equation}
\Lambda_\star-m_1^2\ \ge\ \delta\ >\ 0 ,
\label{eq:est-i2}
\end{equation}
and the same holds in the $2^{++}$ channel with its own $\nu^{(2)}$, $m_2$ and
$\Lambda^{(2)}_\star$. \LT.
\end{claimbox}

\noindent\textbf{Why the uniformity in $L$ is the hard half.} Isolation at one $L$ is a statement
about one transfer matrix; isolation uniformly in $L$ is a statement about a one-parameter family,
and gaps along a family are exactly what degeneracies destroy: the von Neumann--Wigner phenomenon
\cite{vnw} makes a crossing generic once the symmetry that protects it is broken. The five rings earn
the uniformity; an estimate at fixed $L$ would not.

\noindent\textbf{The five rings are the same five.} Isolation is not a second chain: rings R1--R5 of
\S\ref{sec:est-i1} certify \eqref{eq:est-f1} and \eqref{eq:est-i2} on the same transfer matrix, and
the five holes A--E are the same five holes. This is why the two statements share a subsection in
the corpus and are separated here only for readability.
\src{\code{docs/paper-v2/ch20-unconditional-i1.tex}, its five-ring section and its five-hole table;
\code{Y02-CHANNELS.md}:16--20.}

\smallskip
\noindent\textbf{Proposition C, machine-verified.} The isolation has a finite-dimensional shadow that
is a Lean theorem. In the finite-support model $\mathrm{Spec}=(\text{points},\text{weights})$, put
$I2:=\exists\Lambda\,(\Lambda$ is the least support point above $m_1$ and $m_1^2<\Lambda)$ and
$\mathrm{SecExists}:=\exists x\in\text{points},\,m_1<x$. Then
\begin{equation}
0<m_1<1\quad\Longrightarrow\quad\bigl(I2\ \Longleftrightarrow\ \mathrm{SecExists}\bigr),
\label{eq:est-propc}
\end{equation}
with \code{sorry} $=0$. The forward direction is unconditional; the backward direction is exactly
where the window $0<m_1<1$ is spent, through $m_1^2<m_1$.
\src{\code{L04PropC.lean}: \code{propC}, \code{propC_fwd}, \code{propC_bwd}, \code{sq_lt_self};
\code{LOOP-RUNBOOK.md}:9109--9129.}

\smallskip
\noindent\textbf{Two limits of Proposition C, recorded because they are load-bearing.}
\emph{(i)} $I1$ is automatic only when the point table is declared to be the positive-weight support;
with a bare grid of points $I2$ does not imply $I1$, and the file contains that counterexample
(\code{counterexample_I2_not_I1}). \emph{(ii)} The convexity lemma $A1$ is machine-verified in its
finite-sum form only: \code{MeasureTheory} is unavailable in the build, so the integral form stays
\LC{} (\code{L04PropC.lean}, header and \code{a1}; \code{LOOP-RUNBOOK.md}:9129). Neither limit
touches \eqref{eq:est-i2}, which is not a finite-support statement.

% =====================================================================
\subsection{Uniqueness of the target: \texorpdfstring{$N09$-$1$}{N09-1}}
\label{sec:est-n09}

\begin{claimbox}{cC}
\claimhead{Theorem}{(uniqueness of the target, $N09$-$1$; conditional on the M6 specification and
the $L\to\infty$ reading).} Let a quantity $X$ be required to pass the four-model slot test --- free
field $0$, $2$D Yang--Mills ``true'', M6 ``false'', single plaquette ``undefined'' --- with the four
slots kept distinct. Then $X$ is a function of the mass: there is $f$ with $f(0)=0$ and $f>0$ for
$m>0$ such that $X=f(m)$. \LC. Consequently the target is unique: every admissible ``new
invariant'' is a relabelling of $m$.
\end{claimbox}

\noindent\textbf{The grade.} The two hypotheses --- the M6 specification and the $L\to\infty$ reading
of its spectrum --- are printed in the statement, and by the grading rule of the title page a theorem
carrying printed hypotheses is \LC, not \LT. The earlier edition graded this statement \LT; that is
the same species of over-report as the window corrected in \S\ref{sec:est-window-fix}, and it is
corrected here rather than quietly. The content is unaffected: the hypotheses are the ones under
which $M6$ was constructed, and no use of $N09$-$1$ in this paper leaves them.

\noindent\textbf{Why.} M6 is oscillator $\otimes$ free mode with $w=1$, and by the $L\to\infty$
reading its only vanishing scale is the spectral bottom $m_6=\min(w,2\pi^2/L^2)$. A quantity that is
defined on M6 and not identically true must take the value $0$ there, i.e.\ must vanish with
$m_6\to0$; on $2$D Yang--Mills it is positive where $m=3g^2/8>0$. A functional that is defined,
non-vacuous and zero exactly on the degenerate model is therefore a monotone relabelling of $m$ at
$0$. \src{\code{N09-RESTATE.md}:44--53 (\S2, the four numbered steps), :24--30 (four-valued
signature).}

\smallskip
\noindent\textbf{The qualification, and it is the thirty-eighth ledger entry.} Statement
$N09$-$1$ holds \emph{only at fixed operator $O$}. The four-model test fixes the \emph{model}; it
does not fix the \emph{operator}, and the atom weight is a property of the pair. On one ensemble, one
lattice, one $\beta$, with no change of $H$, the same Markov chain and the same $12000$
configurations give
\begin{equation}
\begin{gathered}
w_1=0.3427\pm0.0119\quad(O=\tfrac14\operatorname{tr}F^2,\ \text{unsmeared}),\\
w_1=0.0168\pm0.0142\quad(\text{spatial APE }n=1),
\end{gathered}
\label{eq:est-w1op}
\end{equation}
a ratio $0.343/0.017\approx20.2$. The correct reading of $N09$-$1$ is therefore
$X=f(m;O)$: uniqueness is uniqueness \emph{per operator}, and any statement of the form
``$w_1>0$'' has no content until $O$ is named.
\begin{sloppypar}
\noindent\src{\code{docs/paper-v7/w1/main.tex}:233, :251, :261--264 (the pair and the ratio
$20.2$); \code{position-current.html}:377 (entry 38); \code{ATOM2.md}:15--16 for the one-line
counterexample $w_1=\cos^2\theta$ at fixed spectrum.}
\end{sloppypar}

% =====================================================================
\subsection{The no-go \texorpdfstring{$SG1$-$1$}{SG1-1}: \texorpdfstring{\LC}{[C]}, not \texorpdfstring{\LT}{[T]}}
\label{sec:est-sg1}

\begin{claimbox}{cC}
\claimhead{Theorem}{(no-go $SG1$-$1$; grade \LC).} On the class $\mathcal M$ of non-negative Borel
measures on $[0,\infty)$ with finite moments up to a fixed $N$ and with
$s_1(\nu)=\inf\operatorname{supp}\nu\in\mathbb R$, there is no functional
$F:\mathcal M\to\mathbb R$ satisfying simultaneously
\begin{equation}
\begin{gathered}
\textbf{(i)}\ \text{$G1$-covariant},\qquad
\textbf{(ii)}\ F\le s_1,\qquad
\textbf{(iii)}\ \text{finite-moment computable},\\
\text{hence}\qquad
\textbf{(i)}\wedge\textbf{(ii)}\wedge\textbf{(iii)}\ =\ \varnothing .
\end{gathered}
\label{eq:est-sg1}
\end{equation}
\textbf{The grade is \LC{} and not \LT}, and \S\ref{sec:est-sg1-grade} states exactly which step
costs the downgrade.
\end{claimbox}

\begin{proofx}
Fix $a\ge0$, $\lambda>0$ and $0<\varepsilon<\min(a,2\lambda)$, and put
\begin{equation}
\begin{gathered}
\nu_\varepsilon^{(a)}\ =\ \frac{\varepsilon}{2\lambda}\,\delta_\varepsilon
\ +\ \Bigl(1-\frac{\varepsilon}{2\lambda}\Bigr)\,\delta_a\ \in\ \mathcal M ,
\qquad
s_1\bigl(\nu_\varepsilon^{(a)}\bigr)=\varepsilon ,\\
m_k\bigl(\nu_\varepsilon^{(a)}\bigr)\ =\ \frac{\varepsilon}{2\lambda}\,\varepsilon^{k}
+\Bigl(1-\frac{\varepsilon}{2\lambda}\Bigr)a^{k}
\ \xrightarrow[\varepsilon\downarrow0]{}\ a^k=m_k(\delta_a)
\qquad(k\ge1).
\end{gathered}
\label{eq:est-atom}
\end{equation}
\emph{Step 1.} By \eqref{eq:est-atom} the moment vector of $\nu_\varepsilon^{(a)}$ converges to that
of $\delta_a$. \emph{Step 2.} Continuity of $\Phi$ in (iii) gives
$F(\nu_\varepsilon^{(a)})\to F(\delta_a)$, while (ii) gives
$F(\nu_\varepsilon^{(a)})\le s_1(\nu_\varepsilon^{(a)})=\varepsilon$; hence $F(\delta_a)\le0$ for
every $a\ge0$. \emph{Step 3.} $\tau_a\delta_0=\delta_a$, so (i) gives
$F(\delta_a)=F(\delta_0)+a$; combined with Step 2 this forces $F(\delta_0)+a\le0$ for every $a\ge0$,
which fails for $a>|F(\delta_0)|$. No such $F$ exists. The argument uses one moment ($N\ge1$) and
continuity, and nothing else.
\end{proofx}

\noindent The three-line skeleton is the corpus statement, transcribed word for word: the atom family
drifts the moments by $O(\varepsilon)$ while $s_1:1\to\varepsilon$; the machine check at
$a=0.5/1/2/5$ returns \textsf{NO} in all four rows; hence the only escape is the full-moment limit,
the only reachable object there is $F=s_1$, so $G\equiv0$, which collides with $N09$-$1$.
\src{\code{SG1.md}:62--67 as quoted in \code{docs/paper-v7/sg1/main.tex}:252--268;
\code{w7/sg1_G.log}:37--42 (the four \textsf{NO} rows), reproduced at
\code{docs/paper-v7/sg1/main.tex}:401--404 (values $1.0\cdot10^{-6}$, $2.49999\cdot10^{-7}$,
$4.99999\cdot10^{-7}$, $9.99999\cdot10^{-7}$, $2.5\cdot10^{-6}$ against $s_1$ equal to
$0.5/1/2/5$).}

% ---------------------------------------------------------------------
\subsubsection{Why the grade of \texorpdfstring{$SG1$-$1$}{SG1-1} is \texorpdfstring{\LC}{[C]}: the thirty-seventh ledger entry}
\label{sec:est-sg1-grade}

\noindent Condition (iii) admits exactly one reading in the proof above: \emph{there is an
$N<\infty$, fixed independently of $\nu$, and a continuous $\Phi:\mathbb R^N\to\mathbb R$ with
$F(\nu)=\Phi(m_1(\nu),\dots,m_N(\nu))$ for every $\nu\in\mathcal M$}. Under that reading Step 2 is a
one-line consequence of continuity, and \eqref{eq:est-sg1} is a theorem.

\smallskip
\noindent\textbf{If (iii) is read as computability} --- recursion-theoretic, or ``computable from
$(S_W,\beta,L,a)$ by an algorithm'' --- then the passage
$F(\nu_\varepsilon^{(a)})\to F(\delta_a)$ does not follow: a functional computable in the moment
data but discontinuous in the moment topology is outside the killed class. \textbf{That single step
is therefore \LC, and with it the true grade of $SG1$-$1$ is \LC, not \LT.} The proof is three lines
long only because continuity is doing all of the work; remove continuity and nothing is left.
\src{\code{docs/paper-v7/sg1/main.tex}:207--232 (convention \code{iii-strict} and the remark on
(iii)), :320--330 (Step 2 carries the chip); ledger entry 37,
\code{position-current.html}:378.}

\smallskip
\noindent\textbf{Two conventions that the statement needs, and neither is cosmetic.}
\emph{(1) (i) is pinned to the G1 energy axis.} Under a spatial translate the object
$\inf\operatorname{supp}\nu$ is invariant, so a G2 reading makes (i) constrain nothing and
simultaneously puts it in direct conflict with (ii); the conflict is an artefact of the reading.
\emph{(2) (ii) quantifies over the whole measure class} $\mathcal M$, not merely over the spectral
measures physics supplies. That makes the emptiness \emph{stronger} than its physical-class version:
a functional that fails on the larger class fails a fortiori on the subclass.
\src{\code{B05-DUAL.md}:76--83, :80; \code{docs/paper-v7/sg1/main.tex}:163--205 (conventions
\code{domain} and \code{G1}).}

\smallskip
\noindent\textbf{Relation to $N09$-$1$: no implication in either direction.} $N09$-$1$ quantifies
over \emph{quantities on models} that pass the slot test; $SG1$-$1$ quantifies over
\emph{functionals on measures}. They meet at exactly one point, and that point is a degenerate end:
the survivor $F=s_1$ gives $G\equiv0$, which passes $4/4$ and lands inside $\operatorname{dom}(N09$-$1)$
as the function $f\equiv0$. The corpus phrase ``collides with $N09$-$1$'' therefore means: the
non-trivial defect $G$ is excluded, and the zero defect is not an object at all. It does not mean
that one statement entails the other; the two quantify over different carriers, and each is
conditional on a hypothesis the other does not share --- the $M6$ specification for $N09$-$1$, the
fixed-$N$ continuity reading of (iii) for $SG1$-$1$ --- so no grade transfers between them.
\src{\code{docs/paper-v7/sg1/main.tex}:617--641 (remark \code{n09rel}), :441--448 (signature of the
survivor); \code{N09-RESTATE.md}:44--53.}

% =====================================================================
\subsection{Four machine-verified lemmas: two counterexamples, two certificates}
\label{sec:est-counter}

\noindent These four are model-level or arithmetic results, hence \LM: they are verified in a
solvable model or by exhaustive machine enumeration. Two of them kill implications and two certify
arithmetic, and the distinction matters, because the second pair is not evidence about Yang--Mills at
all (\S\ref{sec:est-window-fix}).

\begin{enumerate}[label=\textbf{\arabic*.},leftmargin=2.0em]
\item \textbf{What it kills:} $G''>0\Rightarrow m>0$ (``$G''$ sufficient''). \emph{Witness:} M6 has
      $G''=1/\mathrm{Var}(a)=2\omega>0$ while its spectral bottom obeys
      $m_6=\min(\omega,2\pi^2/L^2)\to0$; the finite-dimensional $2\times2$ witness is
      machine-checked, \code{sorry} $=0$, axioms
      $[\code{propext},\code{Classical.choice},\code{Quot.sound}]$. The file records that this
      witness realises the failure through a degenerate vacuum, so it kills the bare implication and
      is not claimed as an M6-class counterexample.
      \begin{sloppypar}
      \noindent\src{\code{L05Counterexamples.lean}; \code{L05-COUNTER.md};
      \code{LOOP-RUNBOOK.md}:8691--8697, :8841; \code{N09-RESTATE.md}:66 (the $m_6$ formula).}
      \end{sloppypar}
\item \textbf{What it kills:} $m>0\Rightarrow f_1^2>0$ (``$f_1^2$ necessary''). \emph{Witness:} in
      the $2$D Yang--Mills type model $[H,O]=0$ with a non-degenerate spectrum forces the
      off-diagonal element to vanish, $f_1^2=0$, while the gap stays positive; the explicit
      $2\times2$ witness is $H=\mathrm{diag}(0,1)$, $O=\mathrm{diag}(3,-1)$ with gap $1$.
      \src{\code{M03-COUNTER.lean} (\code{f1sq_eq_zero}); \code{LOOP-RUNBOOK.md}:9683--9731.}
\item \textbf{What it certifies:} the window's arithmetic at the elementary-bound step,
      $\exp(5/96)\le1.054$ and $4\exp(5/96)<7.0633$; the partial sum is $1.0534632$, the remainder
      $3.83\cdot10^{-7}$, the margin $4.6\cdot10^{-4}$.
      \src{\code{LOOP-RUNBOOK.md}:8961--8962; \code{L06-INTERVAL.md}:68--69;
      \code{L02-CBARKATO.md}:84--85.}
\item \textbf{What it certifies:} the window arithmetic as a whole --- twenty-five theorems,
      \code{sorry} $=0$, twelve axiom checks clean, including $8>7.0632853$ and the rational
      certificate. \src{\code{M02-CBAR.lean}; \code{LOOP-RUNBOOK.md}:9512--9528.}
\end{enumerate}

\noindent\textbf{How the four are used, and how they are not.} Rows 1 and 2 kill two implications
that the project had used as inference rules --- ``$G''$ is the key'', ``$f_1^2$ is the key'' --- and
they kill them on models where the implication is checkable by kernel, not by sampling. Rows 3 and 4
certify arithmetic, and arithmetic is not a statement about Yang--Mills: they certify that the numbers
in \S\ref{sec:est-window-fix} are what they are claimed to be. In particular row 4 does not make the
\emph{window} a theorem, because a numerical maximum of an explicit function is still a numerical
maximum; \LM{} is exactly the right grade for it, and it is the grade printed.

% =====================================================================
\subsection{One new \texorpdfstring{\LT}{[T]} this round: the block-Loewner order}
\label{sec:est-loewner}

\begin{claimbox}{cT}
\claimhead{Theorem}{(block-Loewner monotonicity in $\beta$).} Split the transfer kernel as
$T_\beta=D_\beta^{1/2}K_\beta D_\beta^{1/2}$ with $D_\beta=e^{-\beta S_s/2}$ and $K_\beta$ the
half-plaquette kernel. Then for $\beta>\beta'$
\begin{equation}
K_\beta\ \succeq\ K_{\beta'} ,
\label{eq:est-loewner}
\end{equation}
in the Loewner order on $L^2(dU)$: the block coefficients satisfy $c_r(\beta)\ge0$ and are
non-decreasing in $\beta$ for every representation $r$, with $c_r(\beta)=\langle\chi_r\rangle/d_r$, and
the tensor-product multiplicities entering the composition of blocks are non-negative.
\LT.
\end{claimbox}

\begin{proofx}
The order \eqref{eq:est-loewner} reduces, term by term, to $c_r(\beta)\ge c_r(\beta')$ for all $r$:
positive definiteness of $\chi_r(UV^\dagger)$ makes the block comparison a comparison of coefficients
rather than of kernels. Positivity $c_r\ge0$ and monotonicity in $\beta$ follow from the
Griffiths/Bochner route applied to $\langle\chi_r\rangle$ \cite{griffiths,ginibre,ksherman}, and the
composition of two blocks is a sum over $r_1\otimes r_2$ with non-negative multiplicities, so no
negative weight can enter the comparison.
\end{proofx}

\noindent\textbf{What it repairs, precisely.} \code{D05-FKG.md} adjudicated four readings of ``the
order of $T_\beta$'' and left exactly one slot open: at \emph{block} level the coefficient
monotonicity was known for $\mathrm{SU}(2)$ and marked \LN{} for $\mathrm{SU}(N\ge3)$
(\code{D05-FKG.md}:69, :100). \eqref{eq:est-loewner} discharges that slot, so the block-level row of
that adjudication changes from \LN{} to \LT. \src{\code{position-current.html}:380 (the round-178
record of the new \LT, item 6); \code{D05-FKG.md}:62--71.}

\smallskip
\noindent\textbf{What it does not repair.} The \emph{full-kernel} order $T_\beta\succeq T_{\beta'}$
is not established, and \eqref{eq:est-loewner} does not touch it: the factor $D_\beta$ decreases
pointwise in $\beta$, and that obstacle sits in the other slot of the splitting. And even if the
block order is granted, what it transports is the \emph{unnormalized} eigenvalue sequence --- not the
normalized kernel, not the eigenvectors, and not $F$ or $m$. The repair is a repair of one slot, not
a route to the gap. \src{\code{D05-FKG.md}:71 (the $D$-factor obstacle), :76--82 (``output zero for
$F$, for $m$, for $\mathbb G$''), :100 (the row table).}

\begin{warnbox}
\textbf{A bug, caught in the same round, that voids earlier numerical sign determinations.} A
QR-based sampler fabricates ``sign changes of $c_r$'': the sampled coefficients flip sign as a
numerical artefact of the factorization, not as a property of the kernel. Consequences, both
recorded and neither negotiable: (a) four places in the corpus state the coefficient convention in a
form that the artefact made look consistent, and those statements have to be changed; (b) several
earlier numerical conclusions of the form ``$c_r$ changes sign here'' are void and are withdrawn
pending re-derivation. The analytic statement \eqref{eq:est-loewner} is independent of the sampler,
which is why it keeps \LT{} while the numerical readings do not.
\src{\code{position-current.html}:417 (the bug entry), :399 (``two solid items: block-Loewner \LT{}
and the QR-sampler bug'').}
\end{warnbox}

% =====================================================================
\subsubsection*{What this section does not claim}

\begin{enumerate}[label=(\roman*),leftmargin=2.2em]
\item \textbf{Not Clay.} No continuum mass gap is claimed, $a\to0$ is not taken, and no statement
      here crosses to $\beta\approx2.5$. The residue and the isolation are strong-coupling
      statements at fixed spacing, and their window \eqref{eq:est-window2} is a window in the
      coupling, not in the spacing.
\item \textbf{The positivity of the mass inside the window is not ours.} For $\beta$ below the
      strong-coupling radius, $m>0$ is Osterwalder--Seiler's cluster expansion \cite{osterwalder-seiler}
      and the edge $1/48$ of the window is where Shen--Zhu--Zhu's log-Sobolev and Poincar\'e
      inequalities are proved \cite{szz}; the same conclusion in
      the polymer language is \cite{seiler,bgdf}, and for the lower-dimensional models
      \cite{gjs,fo}. What \S\ref{sec:est-i1}--\S\ref{sec:est-i2}
      add is the residue $I1$ and the isolation $I2$, i.e.\ the removal of the premise that the
      $0^{++}$ channel is strictly lightest --- and nothing more.
      \src{\code{docs/paper-v2/ch20-unconditional-i1.tex} \S D.2.}
\item \textbf{The step from \eqref{eq:est-i2} to $m_1>0$ is vacuous and is not used.} The side
      condition of the Kato-type reading is verbatim $m_1>0$, so the implication is written in
      through its own hypothesis; it is recorded as a qualification, never as a gain.
      \src{\code{docs/paper-v2/ch20-unconditional-i1.tex} \S D.1; \code{G07-KATOADJ.md} \S0 F.}
\item \textbf{Proposition C is a finite-dimensional shadow.} It covers a finite support with a fixed
      point table; it does not discharge any statement uniform in $(L,a)$ and must not be read as
      closing a gap in $\beta$. \src{\code{N09-RESTATE.md}:66--70 (\S5 item 4).}
\item \textbf{The block-Loewner order produces no lower bound.} By construction it transports
      eigenvalues, and the project's target is a lower bound on a spectral bottom; the two are
      different objects, and \S\ref{sec:est-loewner} says so in its own scope sentence.
\end{enumerate}

\noindent Everything in this section is therefore of one of two kinds: a statement whose proof is
complete at the grade printed, or a statement whose missing step is named at the place where it is
missing. There is no third kind in the register, and no claim above is graded above its weakest link.

\endgroup

% =====================================================================
%  PART III -- THE MEASUREMENT  (the only positive result in the paper)
%  Replaces the stub.  English only.  Every number carries its source.
%  Primary sources:
%    research/v3/species/RUN1.md      (beta = 2.5, L_s = 6, 8; run1)
%    research/v3/species/RUN3.md      (beta = 2.5, L_s = 8, N = 1.2e4; run3)
%    research/v3/species/RUN2.md      (estimator identity, m_1 window)
%    docs/position-current.html       (cost origin, four-model ledger)
%    docs/paper-v2/ch12-numerics.tex  (beta = 2.2 volume family)
%    docs/paper-v7/routes20/main.tex  (contact term M_0 ~ a^-8)
%    docs/yang-mills-frontier.md      (gamma_{w_1} = 0)
% =====================================================================

\clearpage
\part{The measurement}

% Section numbering runs continuously through the parts, so that no two
% sections of this paper carry the same number.  (The v3 source restarted
% the counter here, which produced a second \S3, \S4, \S5 and \S6.)
% Cross-references use \label/\ref throughout and are unaffected.

\section{Measuring the atom fraction \texorpdfstring{$w_1$}{w1} at \texorpdfstring{$\beta=2.5$}{beta=2.5}}
\label{sec:measurement}

This part carries the only positive result of the paper, and it is a measurement, not a theorem:
everything in it is graded \LE{} except where a statement is arithmetic or definitional. We define
the quantity (\S\ref{sec:m-what}), report what the lattice actually cost (\S\ref{sec:m-setup}),
give the numbers (\S\ref{sec:m-results}), show that they are determined by the operator and not only
by the theory (\S\ref{sec:m-opdep}), list the three systematics that limit them (\S\ref{sec:m-syst}),
record the two bugs found while producing them (\S\ref{sec:m-bugs}), and say in one place what the
measurement does and does not establish (\S\ref{sec:m-scope}). The order is deliberate: the negative
half of this part is longer than the positive half, and the reason is in \S\ref{sec:m-scope}.

\subsection{What is measured}
\label{sec:m-what}

The object is the atom fraction of the spectral measure of the $0^{++}$ channel at zero momentum.
Fix the operator, its connected two-point function and its spectral measure,
\begin{equation}
O=\tfrac14\operatorname{tr}F^2,
\qquad
C(t)=\langle O(t)O(0)\rangle_c=\int e^{-mt}\,d\nu(m^2),
\qquad
M_0=C(0)=\operatorname{Var}(O),
\label{eq:m-defs}
\end{equation}
let $m_1$ be the lowest mass in $\nu$ and $f_1^2=\nu(\{m_1^2\})=|\langle G|O|\Omega\rangle|^2$ the
weight of the lowest atom. The measured quantity is
\begin{equation}
w_1 \;=\; \frac{f_1^2}{M_0} \;=\; \frac{\nu(\{m_1^2\})}{\nu(\mathbb{R}_+)}\in[0,1],
\qquad
\gamma_{w_1}=0 .
\label{eq:m-w1}
\end{equation}

\begin{claimbox}{cE}
\claimhead{Claim}{(definition, and the one identity that makes it a number).}
For every gauge-invariant $O$ with $\|O\|<\infty$ one has $0\le w_1\le1$; and $w_1$ is invariant
under operator rescaling $O\mapsto ZO$.
\end{claimbox}

\begin{proofx}
$f_1^2=\nu(\{m_1^2\})\le\nu(\mathbb{R}_+)=M_0$ is the statement that one atom is at most the total
mass of a positive measure, and positivity gives $f_1^2\ge0$. For the second part, both $f_1^2$ and
$M_0$ are quadratic in $O$, so $O\mapsto ZO$ multiplies numerator and denominator by $|Z|^2$ and the
ratio is unchanged; hence the anomalous dimension of $w_1$ vanishes, $\gamma_{w_1}=0$. This is the
whole content: the identity fixes the value of $w_1$ under a change of normalisation, and fixes
nothing about its sign. Positivity of $w_1$ is not a scaling statement and is not obtained here.
\end{proofx}

The scaling identity is what makes the number meaningful, and its failure at finite cutoff is the
reason a scheme has to be named before any number is quoted:
\begin{itemize}
\item the numerator is an infrared quantity, a point mass at the lowest atom;
\item the denominator is dominated by the ultraviolet contact term of $O^2$ at coincident points,
  $M_0^{\mathrm{bare}}\sim a^{-8}$;
\item therefore $w_1^{\mathrm{bare}}=f_1^2/M_0^{\mathrm{bare}}\to0$ as $a\to0$. The ratio is dragged
  to zero by the denominator, and the vanishing carries no information about the spectrum.
\end{itemize}
So $w_1$ is not a number but a family $w_1(L,\beta,t)$, with $t$ the renormalisation scheme, i.e.\
the prescription that removes the contact term from $M_0$. Every number in this part belongs to one
explicit member of that family, and to exactly one:
\begin{equation}
w_1 \;=\; \frac{A}{C(0)},
\qquad
C(t)=\frac{1}{L_t}\sum_{t_0}\langle S(t_0+t)S(t_0)\rangle_c,
\qquad
S(t)=\sum_{\mathbf{x}}P(\mathbf{x},t),
\label{eq:m-est}
\end{equation}
with $C(t)$ fitted by $A\,(e^{-mt}+e^{-m(L_t-t)})$ on $t=1,2,3$ at fixed $L_t$, the mass pinned at
$m_1=1.546$, and $P=\tfrac13\sum_{i<j\le3}\tfrac12\operatorname{Re}\operatorname{tr}U_{ij}$ the
zero-momentum spatial plaquette sum. In words: the scheme $t$ here is \emph{the bare operator with
the contact term left in the denominator}, and the measured value is the ratio of the fitted
amplitude to the unsubtracted $t=0$ correlator of the same operator at the same lattice spacing. It
is the only scheme the present data define, and it is the scheme to which $0.3427$ belongs. In
particular it is a finite-$a$ statement: nothing below is a continuum value of $w_1$, and
\S\ref{sec:m-scope} says so again in the negative list.
\src{Definitions and the estimator: \code{eq:m-defs}--\code{eq:m-est} follow the pinned conventions
of \code{RUN1.md:31--39} and \code{RUN3.md:41--47} verbatim (operator, zero-momentum sum, periodic
time averaging, fit window $t=1,2,3$ at both volumes, $w_1=A/C(0)$, $20$-block delete-one jackknife
errors). The pinning mass $m_1=1.546$ is the lower edge of the allowed window derived in
\code{RUN2.md:101--104}. $\gamma_{w_1}=0$ with the reading ``$f_1^2$ and $M_0$ carry the same
$Z_O^2$'' is \code{docs/yang-mills-frontier.md:339}. The contact term
$M_0^{\mathrm{bare}}\sim a^{-8}$ and the resulting $w_1^{\mathrm{bare}}\to0$ are
\code{docs/paper-v7/routes20/main.tex:168--177}.}

\subsection{Lattice setup and cost: how eleven hours became 265 seconds}
\label{sec:m-setup}

The run is four-dimensional $\mathrm{SU}(2)$ Wilson with $\beta=2.5$ ($\kappa=\beta/2=1.25$), spatial
volumes $L_s=6$ and $L_s=8$, temporal extent $L_t=12$, periodic boundary conditions, a
Kennedy--Pendleton heat bath compiled with \code{numba.njit} and \code{fastmath}, and sequential
single-chain updates. The observable is the zero-momentum spatial plaquette sum \eqref{eq:m-est};
one configuration yields the correlator at every operator level at once.

The estimate this run replaces was eleven hours per core for a single volume. The cause was not the
machine and not the algorithm: the plan thermalised each of $10^{5}$ independent configurations
separately, that is, it paid $10^{5}$ times for $400$ equilibration sweeps each. The fix is one line
of logic --- thermalise once per chain, then produce configurations consecutively along that chain,
and measure all operators on each configuration --- and it changes the cost by two orders of
magnitude. Measured:

\begin{table}[htbp]
\centering\small
\begin{tabular}{@{}llrrrr@{}}
\toprule
volume & lattice (sites / links) & $N$ & sweeps/cfg & ms/sweep & wall clock \\
\midrule
$L_s=6$  & $6^3\times12$ \ (2592 / 10368) & $6000$  & $6$ & $2.223$ & $290$ s \\
$L_s=8$  & $8^3\times12$ \ (6144 / 24576) & $4152$  & $3$ & $5.390$ & $265$ s \\
$L_s=8$  & $8^3\times12$, eight chains $\times\,1500$ & $12000$ & $3$ & --- & $34.0$ s \\
\bottomrule
\end{tabular}
\caption{Cost of the measurement. The first two rows are single-threaded, no parallel; the third is
the same volume with eight independent chains over \code{prange}, $352.6$ configurations per second.
Sources: \code{RUN1.md:16--25} (rows 1--2), \code{RUN3.md:20--37} (row 3).}
\label{tab:m-cost}
\end{table}

The correctness check is the plaquette,
$\langle\tfrac12\operatorname{Re}\operatorname{tr}U_p\rangle=0.65270\pm0.00005$ at $L_s=6$,
$0.65236\pm0.00004$ at $L_s=8$, and $0.65234\pm0.00002$ in the independent eight-chain run, all
three consistent with the values tabulated for $\mathrm{SU}(2)$ in \cite{creutz} ($\approx0.65$ at
this coupling). The sampler is therefore
right; the estimator is a separate matter, and \S\ref{sec:m-syst} is about the estimator.

\begin{qbox}
\textbf{The eleven-hour estimate was an implementation problem, not a machine problem.} Eleven hours
per core against $265$ s measured is a factor $1.5\times10^{2}$, and every digit of that factor is
accounted for by re-thermalisation, not by hardware. The same correction is what makes the
statistics affordable later: the run that \S\ref{sec:m-syst} says is still needed
($N\approx4.8\times10^{4}$) costs about $136$ s at the measured throughput, not eleven hours.
\end{qbox}
\src{Cost origin and the ``$10^{5}$ thermalisations'' diagnosis: \code{docs/position-current.html:406}
(``the $11$ h/core of a1 comes from thermalisation being repeated $10^{5}$ times''). Plaquettes and
sweep times: \code{RUN1.md:23--25}, \code{RUN3.md:20--21,33}. The factor
$1.5\times10^{2}=11\,\mathrm{h}/265\,\mathrm{s}$ is arithmetic on these two numbers.}

\subsection{Results}
\label{sec:m-results}

\begin{table}[htbp]
\centering\small
\begin{tabular}{@{}llrrrr@{}}
\toprule
operator & volume & $w_1$ (at $m_1=1.546$) & $w_1/\sigma$ & $N$ & run \\
\midrule
raw ($\tfrac14\operatorname{tr}F^2$) & $L_s=6$, $L_t=12$ & $0.373\pm0.018$  & $20.7\sigma$ & $6000$  & run1 \\
raw ($\tfrac14\operatorname{tr}F^2$) & $L_s=8$, $L_t=12$ & $0.342\pm0.015$  & $22.8\sigma$ & $4152$  & run1 \\
raw ($\tfrac14\operatorname{tr}F^2$) & $L_s=8$, $L_t=12$ & $0.3427\pm0.0119$ & $28.8\sigma$ & $1.2\times10^{4}$ & run3 \\
\bottomrule
\end{tabular}
\caption{The measurement. Errors are $20$-block delete-one jackknife. Sources:
\code{RUN1.md:47,52} (rows 1--2), \code{RUN3.md:51} (row 3); the $w_1/\sigma$ column is the ratio of
the two columns to its left ($0.373/0.018=20.7$, $0.342/0.015=22.8$, $0.3427/0.0119=28.8$).}
\label{tab:m-results}
\end{table}

Three readings of Table~\ref{tab:m-results}, in increasing order of how much they cost:

\begin{enumerate}
\item \emph{Reproducibility.} The two $L_s=8$ runs are independent: different seeds, $8$ chains
  instead of one, $N$ larger by a factor $2.9$. They agree to $+0.0004$, that is $0.02\sigma$. The
  number did not move when the statistics tripled, and that is the strongest single statement in
  this part.
\item \emph{Significance at fixed mass.} $0.3427/0.0119=28.8\sigma$ away from zero. This is the
  significance of $w_1$ with $m_1$ pinned; it is \emph{not} the signal-to-noise of the correlator,
  which is $5.0$ and is discussed in \S\ref{sec:m-syst}. The two numbers are not interchangeable and
  we do not present the larger one as if it were the smaller one's replacement.
\item \emph{Volume dependence.} $(0.373-0.342)/0.373=8.3\%$ between $L_s=6$ and $L_s=8$: the
  fraction does not dilute with volume, in contrast with the free-field behaviour.
\end{enumerate}

The volume trend is available at a second coupling, where the same estimator was run on three
volumes: at $\beta=2.2$ the unsmeared ratio runs $C_1/C_0=0.1367\to0.1229\to0.1216$ for $L=4,6,8$,
a drop and then a plateau rather than a $1/L$ decay, and the three-volume fit of the upper bound
$\hat w_1(t)=e^{mt}C(t)/C(0)$ to $w_\infty+c/L$ gives $w_\infty\approx0.54$. The worst point of that
family is therefore at $L=4$, the smallest volume, and not at $L=\infty$: the finite volume pushes
the fraction \emph{down}, so the $\beta=2.5$ pair sits inside a plateau rather than on a falling
edge. This is a trend, not a limit, and it is the reason we do not extrapolate
$w_1(2.5)$ to $L=\infty$ anywhere in this paper.
\src{Agreement $0.02\sigma$ and the $28.8\sigma$ ratio: \code{RUN3.md:11--12,51}. The $8.3\%$:
\code{RUN1.md:59} and \code{docs/paper-v2/ch12-numerics.tex:281}. The $\beta=2.2$ family
$0.1367\to0.1229\to0.1216$ for $L=4,6,8$, the fit to $w_\infty+c/L$ with $w_\infty\approx0.54$, and
the reading ``inside the plateau, not on a falling edge'': \code{docs/paper-v2/ch12-numerics.tex:284--288};
the same family and the statement that the infimum sits at $L=4$: \code{docs/position-current.html:382}.}

\subsection{Operator dependence}
\label{sec:m-opdep}

This subsection is the most important one in the paper, because it converts the number above from a
fact about Yang--Mills into a fact about a pair.

\begin{table}[htbp]
\centering\small
\setlength{\tabcolsep}{5pt}
\begin{tabular}{@{}lrrrr@{}}
\toprule
basis (same chain, same $12000$ configurations) & $C(0)$ & $C(1)/C(0)$ & $w_1$ & $S/N(2)$ \\
\midrule
raw ($\tfrac14\operatorname{tr}F^2$)      & $19.888$ & $0.07372\pm0.00265$ & $0.3427\pm0.0119$ & $5.0$ \\
spatial APE, $n=1$                        & $45.527$ & $0.00302\pm0.00297$ & $0.0168\pm0.0142$ & $0.9$ \\
spatial APE, $n=2$                        & $39.428$ & $0.00150\pm0.00298$ & $0.0107\pm0.0136$ & $1.6$ \\
\bottomrule
\end{tabular}
\caption{One Markov chain, one set of $12000$ configurations, three operators, all at $m_1=1.546$.
Only the operator changes. Source: \code{RUN3.md:49--53}.}
\label{tab:m-bases}
\end{table}

The ratio of the first row to the second is $0.343/0.017\approx20.2$. Same chain, same
configurations, same jackknife blocks, same fit window: nothing statistical distinguishes the two
numbers, and the difference is not a bug (\S\ref{sec:m-bugs} lists the bugs that were found, and this
is not one of them).

\begin{v3box}
\textbf{The twenty-fold difference is a theorem about notation, realized on data.} Take a
Hamiltonian $H$ whose lowest atom $m_1^2$ is two-fold degenerate, with orthonormal vectors
$\psi_1,\psi_2$, and fix the source family
\[
O_\theta=\mathbf 1+\varepsilon\,(\cos\theta\,\psi_1+\sin\theta\,\psi_2),
\qquad \varepsilon\ll1 .
\]
$H$ is untouched, so the spectrum is untouched: $m_1$, every $m_n$, and the support of $\nu$ are the
same for every $\theta$. But the overlap of the excited part with the lowest atom is
$\varepsilon\cos\theta$, so $f_1^2=\varepsilon^2\cos^2\theta$, while in the scheme $t$ that removes
the contact term carried by $\mathbf 1$ one has $M_0=\varepsilon^2$. Hence
\[
w_1=\cos^2\theta\in[0,1],
\qquad \text{spectrum fixed.}
\]
As $\theta$ runs over $[0,\pi/2]$, $w_1$ sweeps the whole unit interval at a fixed, massive,
unchanged spectrum. The statement is \LT{} --- the proof is \S\ref{sec:reduction-operator} --- and on
these configurations it is additionally \LE, since the factor $20.2$ of Table~\ref{tab:m-bases} is
its instance at $\cos^2\theta=0.049$.
\end{v3box}

The consequence is a statement about language, and it is not negotiable: the sentence
``$w_1(2.5)>0$'' has no content until $O$ is specified. $w_1$ is a property of the pair $(H,O)$ and
not of $H$ alone; Table~\ref{tab:m-bases} is that sentence instantiated on real configurations, with
the factor $20.2$ standing in for the free parameter $\theta$. Two corollaries follow.
\begin{itemize}
\item A change of operator basis is not a change of statistics, so no amount of $N$ will move
  $0.0168$ towards $0.3427$; and no amount of $N$ will make the two agree.
\item The uniqueness statement \texttt{N09-1} --- ``any quantity that passes the four-model test is a
  function of $m$'' --- is true only under \textbf{fixed $O$}. Under a change of $O$ at fixed $H$ the
  quantity moves by a factor $20$ while $m$ does not move at all. This is ledger item $38$, and it is
  the reason the qualifier is now written into the statement rather than left implicit. \LC{}
\end{itemize}
\src{Table~\ref{tab:m-bases} and the factor $20.2$: \code{RUN3.md:49--53}; the same factor quoted as
``$0.343$ vs $0.017$, factor $20$'' at \code{docs/position-current.html:378}. The
$\cos^2\theta$ witness is the one-line counterexample of route \code{u11}, quoted in the abstract of
this paper; the ledger entry is item $38$ at
\code{docs/paper-v7/ledger/main.tex:268}, with the fixed-$O$ qualifier at
\code{docs/position-current.html:378}.}

\subsection{Three systematics}
\label{sec:m-syst}

None of the three is fatal and all three are load-bearing for \S\ref{sec:m-scope}; we state them in
the order in which they would mislead a reader who stopped after Table~\ref{tab:m-results}.

\begin{warnbox}
\begin{enumerate}[label=\textbf{\arabic*.},leftmargin=2.1em]
\item \textbf{The signal-to-noise is below the campaign's own threshold, and the $1/\sqrt N$ route to
  it is dead.} $S/N(2)=C(2)/\sigma_{\mathrm{JK}}=5.0$, with
  $C(2)/C(0)=0.01189\pm0.00238$. The extrapolation used in the planning stage --- $\sigma\propto
  N^{-1/2}$, so $5.9$ at $N=4152$ should become $10$ at $N=1.2\times10^{4}$ --- failed: $N$ rose by
  $2.9\times$ and $S/N$ went from $5.9$ to $5.0$. The requirement is therefore
  $N\approx1.2\times10^{4}\times(10/5.0)^{2}\approx4.8\times10^{4}$, about $136$ s at the measured
  $352.6$ configurations per second. Statistics can be bought; they were not bought here, and the
  threshold $S/N\ge10$ is not met.
\item \textbf{The time correlation is killed by the operator basis, so the GEVP route is unavailable.}
  With spatial APE ($n=1$, $\alpha=0.5$, time direction not smeared, staple normalised $/4$),
  $C(1)/C(0)=0.0030\pm0.0030$ and $C(2)/C(0)=+0.0030\pm0.0032$; $n=2$ gives
  $0.0015\pm0.0030$ and $+0.0034\pm0.0022$. The earlier $/6$ normalisation of run1 gives the same
  $10^{-3}$. So $C(1)\to0$ with $S/N\le1.6$: the projected correlator is a contact term, not a
  propagating state, and a generalised eigenvalue extraction of $m_1$ --- the standard variational
  spectroscopy of \cite{morningstar} --- hence any statement about
  the $\beta$- or $a$-dependence of $w_1$, cannot be built on this basis at $L\le8$.
\item \textbf{The dominant uncertainty is the smearing level, not the volume.} The spread induced by
  the smearing parameter $s$ is $0.29$, against $0.07$ for the volume $L$: a factor $4$. Since the
  $L$-dependence of Table~\ref{tab:m-results} is $8.3\%$, quoting it as an error bar while holding the
  operator fixed would understate the operator dependence of \S\ref{sec:m-opdep} by the very factor
  the quantity is most sensitive to.
\end{enumerate}
\end{warnbox}
\src{(1) $S/N(2)=5.0$, $C(2)/C(0)=0.01189\pm0.00238$, the failed $1/\sqrt N$ extrapolation
($5.9\to5.0$ at $2.9\times$ the statistics), and $N\approx4.8\times10^{4}$:
\code{RUN3.md:13--14,51,84}; the $136$ s is $4.8\times10^{4}/352.6$ on the throughput of
\code{RUN3.md:37}. Note one internal inconsistency in that file, which we resolve in favour of the
table and flag rather than repair: the prose at \code{RUN3.md:14} writes $\pm0.0038$, which would
give $S/N(2)=3.1$ and contradicts its own quoted $5.0$; the table value $\pm0.00238$ reproduces
$0.01189/0.00238=5.0$ and is the one used here. (2) The three APE levels, the $/4$ versus $/6$
normalisation, $S/N(2)\le1.6$ and ``this GEVP route is unusable'': \code{RUN3.md:8--10,52--53,56,66}.
(3) ``$s$-broadening $0.29$ is $4\times$ the $L$-broadening $0.07$'': \code{U14.md:88--89} as
recorded at \code{RUN2.md:8}.}

\subsection{Two bugs found}
\label{sec:m-bugs}

Both are recorded because the second one is the mechanism by which a wrong number would have entered
the paper, and the first one is the mechanism by which a whole class of earlier numbers may already
have.

\begin{qbox}
\textbf{Bug 1: smearing mutated the Markov chain in place.} The first version of the smeared run
passed \code{Ut = U[c]} --- a \code{numpy} \emph{view}, not a copy --- into the APE routine, which
modified it in place. The chain was then continued from the smeared configuration instead of the
sampled one. Symptoms: heat-bath acceptance degraded as $x\to\pm1$, and a single job went from about
$40$ s to more than $10$ min. Fix: allocate \code{buf = np.empty\_like(U[c])} and copy per
configuration. Consequence: \textbf{any earlier result of the form ``smear, then re-measure'' may be
contaminated}, and we do not certify a smeared level produced before the fix. The raw ($n=0$) numbers
of Tables~\ref{tab:m-results} and~\ref{tab:m-bases} are unaffected, since the raw path applies no
smearing at all.

\textbf{Bug 2: the self-contradiction that was not one.} The run chain reported $w_1=0.577$ against
the strict ceiling $\hat w_1(1)=e^{m_1}C(1)/C(0)=0.551$. Since $w_1\le\hat w_1(1)$ is an identity,
the pair $(0.577,\,0.551)$ is impossible, and it was flagged as an internal contradiction. It is not:
the mass had been pinned at $m=1.50$, and the reported value is consistent only for $m_1\ge1.546$.
At $m_1=1.546$ the same data give $0.373\pm0.018$. The resolution is therefore that \emph{the pinned
value was wrong}, and the identity is exactly what caught it. Two consequences we keep: every $w_1$
in this part is quoted at $m_1=1.546$, the lower edge of the allowed window, so they are
\emph{lower edges and not central values}; and ``report $\hat w_1(t)$ alongside $w_1$'' is now a
mandatory step of the analysis, not a courtesy.
\end{qbox}
\src{Bug 1: \code{RUN3.md:76--78}. Bug 2: the ceiling $\hat w_1(1)=e^{m_1}C_1/C_0=0.551$, the
reported $0.577$, the window $m_1\ge1.546$ and the instructive reading of the pair:
\code{RUN2.md:101--104}; the pinned-$1.546$ values are those of Table~\ref{tab:m-results}.}

\subsection{What this does and does not establish}
\label{sec:m-scope}

\begin{v3box}
\textbf{What can be said.} For the raw operator $O=\tfrac14\operatorname{tr}F^2$, at $\beta=2.5$,
$L_s=8$, $L_t=12$, in the scheme \eqref{eq:m-est}, the atom fraction is
\[
w_1=0.3427\pm0.0119 ,
\]
which is $28.8\sigma$ away from zero and reproduced by an independent run at $0.02\sigma$; the
dependence on the volume between $L_s=6$ and $L_s=8$ is $8.3\%$; and at $\beta=2.2$ the family over
$L=4,6,8$ decreases and then flattens, with the infimum at the smallest volume. That is the positive
content of this paper, and it is a measurement (\LE{}), not a floor.
\end{v3box}

\begin{warnbox}
\textbf{What cannot be said.} Four sentences, each of which the numbers above forbid.
\begin{enumerate}[label=\textbf{(\roman*)},leftmargin=2.3em]
\item \emph{``$w_1>0$ for any operator.''} A change of basis on the same configurations gives
  $0.0168\pm0.0142$, a factor $20.2$ smaller (\S\ref{sec:m-opdep}). $w_1$ is a property of $(H,O)$.
\item \emph{``The signal-to-noise is enough.''} $S/N(2)=5.0<10$ (\S\ref{sec:m-syst}, item 1). The
  $28.8\sigma$ of Table~\ref{tab:m-results} is the significance of $w_1$ at pinned mass and does not
  replace it.
\item \emph{``Time correlations are present.''} $C(1)/C(0)=0.003$ and the GEVP route is unavailable
  on these volumes (\S\ref{sec:m-syst}, item 2).
\item \emph{``This advances Clay.''} $w_1$ and $m$ do not imply one another in either direction:
  \begin{itemize}
  \item \textbf{2D Yang--Mills:} $m=3g^2/8>0$ while $\nu_{\frac14\operatorname{tr}F^2}\equiv0$; and
    more strongly, $m(L)=\sigma L\to\infty$ while $w_1=0$. A local operator is blind to the gap.
  \item \textbf{M6} (oscillator $\otimes$ free module): $w_1=0.5$ while $m=0$.
  \item \textbf{Free field:} $w_1=0$ while $m=0$.
  \end{itemize}
  So this measurement does \textbf{not} prove the mass gap, and the direction of the failure is not
  the direction a reader might expect: the failure is not that the number is small or noisy, it is
  that the number is orthogonal to the gap.
\end{enumerate}
\end{warnbox}

One qualification about the fourth item, because it is easily confused with \S\ref{sec:m-opdep}, and
the two must be kept apart rather than merged. \S\ref{sec:m-opdep} varies $O$ at \emph{fixed} $H$:
one theory, two operators, a factor $20.2$. Item (iv) compares \emph{different} $H$: four models, in
each of which $w_1$ and $m$ point in opposite directions. The first says the quantity is not a
function of the theory alone; the second says that even after the operator is fixed, the quantity is
not a function of the gap. Both are needed, and neither implies the other.

\medskip
\noindent\textbf{What the residual actually is.} Not statistics: \S\ref{sec:m-syst} item 1 says the
statistics cost $136$ s and can be bought. Not the volume: item 3 says the volume contributes $0.07$
against $0.29$ for the operator. What is missing is an operator basis that keeps a propagating
correlation at these volumes (item 2), and --- above that --- an argument that connects the atom
fraction to the mass at all, which item (iv) shows cannot be an implication in either direction. The
measurement is honest, reproducible, and it does not move the verdict.
\src{Item (iv) in full, including $m=3g^2/8>0$ with $\nu_{\frac14\operatorname{tr}F^2}\equiv0$, the
stronger $m(L)=\sigma L\to\infty$ with $w_1=0$, M6 with $w_1=0.5$ and $m=0$, and the free field with
$w_1=0$ and $m=0$: \code{docs/position-current.html:392}. The cost claim is
Table~\ref{tab:m-cost} and the $136$ s arithmetic of \S\ref{sec:m-syst}; the $0.07$ versus $0.29$
split is item 3 of \S\ref{sec:m-syst}. The grading \LE{} for the measurement and the refusal to call
it a floor follow the provenance rule of the front matter.}

% =====================================================================
%  SECTION 4 -- the negative survey
%  Twenty routes to uniformity in (L,a): one direction, three layers.
% =====================================================================
\clearpage
\part{The negative survey}
\section{Twenty routes, one direction}\label{sec:survey}

This section is the audit of what was tried and did not work. It is written to a single
standard: every route is stated as the theorem it actually proves, then killed by the
hypothesis that is false, and the failure is attributed to a file. The organisation is
forced by the material. There is one target (\S\ref{sec:survey-target}), one table of
twenty routes plus two by-products (\S\ref{sec:survey-table}), and one obstruction with
three mutually independent layers (\S\ref{sec:survey-three}). What survives is two items,
one of them a bug (\S\ref{sec:survey-survives}).

\subsection{The target}\label{sec:survey-target}

\begin{claimbox}{cN}
\claimhead{The target statement}{.} The uniformity asked for throughout the campaign is the
single sentence
\[
\boxed{\;\inf_{L,t}\, w_1\!\left(L,\,2.5,\,t\right) \;>\; 0\;}
\]
where $w_1=f_1^2/M_0$ is the atom fraction of the K\"all\'en--Lehmann measure
($f_1^2=\nu(\{m_1^2\})$, $M_0=\int d\nu=C(0)=\mathrm{Var}(O)$) for the fixed operator
$O=\tfrac14\mathrm{tr}\,F^2$. This is the gap of \cite{jaffe} read on the atom route of
\S\ref{sec:reduction-equiv}; the reading is not a weakening of the Clay statement, and
\S\ref{sec:survey-three} states exactly how far it is from it.
\src{\code{LOOP-RUNBOOK.md}:18982--18983 (\S160.1); definitions \code{LOOP-RUNBOOK.md}:18950--18956.}
\end{claimbox}

Three bookkeeping facts make this a well-posed sentence, and each was a separate round of
work to establish.

\begin{itemize}
\item \textbf{The lattice spacing is not a free parameter.} At the physical coupling
  $\beta=2.5$ the spacing is pinned, $a=a(\beta)$; the family is a graph, not a rectangle.
  Leaving $a$ free loses $\beta=2.5$; pinning it degrades ``$a$-uniformity'' to a
  statement about one point. Both readings are recorded in the corpus, and the second is
  the one used here.
\item \textbf{The ratio is scale invariant.} The engineering dimension of $w_1$ vanishes
  identically, $\gamma_{w_1}=2\gamma_O-2\gamma_O\equiv 0$, because numerator and
  denominator carry the same anomalous dimension. This is why no renormalisation of the
  bare ratio can move the target; the bare ratio nevertheless drifts to zero, $w_1^{\rm
  bare}\to0$, because the contact term makes $M_0\sim a^{-8}$. That drift is an artefact
  of not fixing the renormalisation scheme, not information about the atom.
\item \textbf{The order of the infima does not matter.} On a product set the three
  infima $\inf_L\inf_t$, $\inf_t\inf_L$ and $\inf_{(L,t)}$ coincide (a two-line lemma).
  ``The order is the difficulty'' is therefore false; what is genuinely open is the
  domain constraint, whether $\lim$ sits inside or outside the $\inf$, and the missing
  $t$-axis.
\end{itemize}

\begin{warnbox}
\textbf{The operator is not decoration.} $w_1$ depends on $O$, so \emph{every} infimum
above is taken at fixed $O$. The witness is one line and exact: on a fixed $H$ take the
legitimate sources $O_\theta=1+\varepsilon(\cos\theta\,\psi_1+\sin\theta\,\psi_2)$; then
$w_1=\cos^2\theta\in[0,1]$ while $\mathrm{Spec}(H)$ does not move at all. Hence $w_1$ is a
property of the pair $(H,O)$, never of $H$ alone, and a statement of the form ``$w_1>0$''
has no content until $O$ is named. This is the same fact that makes the measured factor
$20.2$ between two operators on one Markov chain a feature rather than an error.
\src{\code{U11.md} \S\S1--2, \code{LOOP-RUNBOOK.md}:19241--19252 and :19283--19289 (\S165).}
\end{warnbox}

\subsection{A table of twenty routes, one direction}\label{sec:survey-table}

Twenty routes were attacked, in the sense that each was pushed until it either produced a
uniform lower bound on $w_1$ or produced the exact reason it cannot. All twenty fail, and
they fail in one direction: nine give only upper bounds, nine are circular, two have
explicit counterexamples. Two further items were produced by the same sweep and are kept
because they are usable.

{\footnotesize
\setlength{\tabcolsep}{4pt}
\renewcommand{\arraystretch}{1.12}
\begin{longtable}{@{}>{\raggedright\arraybackslash}p{0.5cm} >{\raggedright\arraybackslash}p{3.15cm} >{\raggedright\arraybackslash}p{4.9cm} >{\raggedright\arraybackslash}p{1.15cm} >{\raggedright\arraybackslash}p{5.1cm}@{}}
\caption{Twenty routes to an $(L,a)$-uniform floor on $w_1$, grouped by the direction of what they
produce, with two by-products. The group headings partition the outputs, not the methods
(\S\ref{sec:survey-table}, notes).}\label{tab:routes}\\
\toprule
\textbf{\#} & \textbf{Route} & \textbf{What it gives} & \textbf{Grade} & \textbf{Cause of death}\\
\midrule
\endfirsthead
\toprule
\textbf{\#} & \textbf{Route} & \textbf{What it gives} & \textbf{Grade} & \textbf{Cause of death}\\
\midrule
\endhead
\bottomrule
\endlastfoot
\multicolumn{5}{@{}l}{\emph{Only upper bounds (9)}}\\
1 & compactness & Arzel\`a--Ascoli transports, does not produce: the output is the identity $\inf_K w_1=w_1(x_0)$ & \LN & $K$ is not compact ($a=0$ and $L=\infty$ have no member); relative compactness does not exclude $\ell=0$\\
2 & reflection positivity & $\widehat w_1(T)\downarrow\ \ge w_1$ \cite{os1973}, i.e.\ an upper bound; the direction is reversed & \LN & RP fixes the sign of the monotone quantity, only constrains $w_1$ from above; $2/(L-1)\searrow0$ and $w_1\equiv0$ both satisfy it\\
3 & Berezin--Lieb & $\sum_j\lambda_j\le\sum_j\varepsilon_j$ is a sum bound: $\lambda_1\le(1/N)\sum\lambda$ & \LN & a trace inequality bounds the bottom from above; its only lower-shaped corollary needs $u\ge0$, and on YM the hypothesis \emph{is} the target\\
4 & probabilistic reading \cite{chatterjee} & $w_1=P(S=m_1^2)$ and $|EX|^2\le E|X|^2$ is an upper bound & \LN & the modulus square is not the expectation of an indicator; $|\langle\mathbb{G}|\phi\rangle|^2$ is a difference of two positive currents, sign indefinite\\
5 & Krein--Rutman residual $D_\infty$ & information content a factor $2.0975$ (window $[0.690,1.448]$) & \LN & fails the free-field acid test: at $\beta=0$, $D_\infty=1$ while $w_1=0$\\
6 & concavity & the chord lemma is true: $f$ concave, $f\ge0$, left end $r_0$ $\Rightarrow$ $f(2.5)\ge r_0(\beta_2-2.5)/(\beta_2-\beta_1)$ & \LN & $w_1$ is not concave: the Lorentzian witness $v^2/(4v^2+D^2)$ changes the sign of $f''$ at $2v/\sqrt3$; ray concavity would force $w_1(\infty)\ge r_0$ against the free end\\
7 & RG modulus $\delta_A\cdot a$ & a pure number, $\delta_A\cdot a=0.5686159$ & \LN & a $0.24\%$ coincidence with $a\Lambda(1/48)=0.5700$ (other $\beta$ give $1.208$, $7.4\times10^3$, $5.2\times10^{10}$); the free field is equally ``constant'' ($2\pi/8=0.7854$)\\
8 & $\Phi$ (object mismatch) & an activity sum $\Phi(0)=\sum_{r\ne0}c_r\dim r=\beta=2N^2u$ & \LN & $\Phi$ is a sum over $f$ and $\bar f$ (each carrying $\dim f=N$); $\bar C$ wants one slab, and one plaquette may take $f$ or $\bar f$ but not both\\
9 & $\varphi_{\rm dec}$ & the decay density of the decoration gas, $\log Z_{\rm dec}=\varphi_{\rm dec}T+\log D_\infty+\varepsilon_T$ & \LN & Fekete already gives its rate and its existence without any positivity; a tool whose output is already a theorem adds nothing\\
\multicolumn{5}{@{}l}{\emph{Circular (9)}}\\
10 & positivity improvement & Perron--Frobenius locks the vacuum at every $\beta>0$ & \LN & assumes $r>0$: a uniform cone-improvement constant is equivalent to a uniform gap, so the hypothesis is the conclusion\\
11 & minimality & $\lambda_1(\nu)\ge\operatorname{ess\,inf}G''$ & \LN & reduces to $\operatorname{ess\,inf}G''=H1$, and the Savvidy tail \cite{savvidy,no,leutwyler} makes $\operatorname{ess\,inf}G''<0$\\
12 & operator inequalities & Temple \cite{temple}, Weyl, Lieb--Thirring, Cheeger, ELS & \LN & $w_1$ is not a function of $\mathrm{Spec}(H)$: at fixed spectrum $w_1=\cos^2\theta\in[0,1]$; the input is the target\\
13 & definability & $o$-minimality can be had on $L/a\le C$ & \LN & the non-circular version requires $(N=\infty,\beta=\infty)$ to be in the family with value $>0$, which is the Clay problem itself\\
14 & strong--weak bridge & $\beta<3.777\times10^{-3}$ gives $f_1^2\ge r_0$; $\beta\to\infty$ gives $m_{\rm lat}=D_0e^{-c\beta}$ & \LN & the two ends are of different types, $662\times$ apart; no interpolation theorem (non-negative real-analytic can be zero in between)\\
15 & $\Delta_0$ & $\exists\Delta_0>0\ \forall a<a_0:\ m_{0^{++}}(a)\ge\Delta_0$ & \LN & M6: $\Delta_{\rm dyn}>0$ while $m=0$ (not sufficient); 2D: $\Delta_{\rm dyn}=0$ while $m>0$ (not necessary); and the floor is the target renamed (\S\ref{sec:reduction-legs}(iii)), which is why the row sits in this group rather than among the counterexamples\\
16 & $T4$ & the same statement in physical units, $=H1$ dimension-corrected & \LN & its antecedent is its conclusion in another notation: $C>0\iff(T4)\iff H1$, so the step consumes what it proves\\
17 & $H1$ & a one-sided variance bound at typical field and $k\to0$ & \LN & $H1\not\Rightarrow m>0$: the thin atom witnesses a gap-free state satisfying $H1$\\
18 & $M(\beta)=m/\Lambda$ & an identity ($\gamma_m=0$): flow invariant and positive by construction & \LN & an identity cannot floor anything; its positivity is free and its uniformity is the assumption\\
\multicolumn{5}{@{}l}{\emph{Counterexamples (2)}}\\
19 & $w_1(L)$ monotonicity & the mechanism family realises both directions & \LN & $w_1=2/(L-1)$ is non-increasing with limit $0$; non-increasing does not imply $\inf>0$; identity of the mechanism family, not a defect of one member\\
20 & tight family $+$ limit atom & a compact family in $\mathcal P([0,3])$, weakly closed & \LN & $\int s\,d\mu\equiv2.4$, $E$ jumps $0.3\to0$, $\mathrm{TV}\equiv0.6$: tightness does not give a uniform floor, the atom survives only at the limit point\\
\multicolumn{5}{@{}l}{\emph{By-products (2)}}\\
21 & block-Loewner & $K_\beta\succeq K_{\beta'}$ for all $N$ & \LT & not a route to $w_1$ (blocked by the $D$-factor and by not transporting eigenvectors) but it repairs the previously unproved step of D05\\
22 & QR-sampler bug & the sampler fabricates sign changes & \LE & four corpus conventions need revision; earlier numerical ``$c_r$ changes sign'' conclusions are void\\
\end{longtable}
}

\src{Row sources, in the order of the table.
1: \code{U01.md} \S\S1--3 (\code{LOOP-RUNBOOK.md}:19064--19070).
2: \code{U03.md} (\code{LOOP-RUNBOOK.md}:19307--19311); the bound $\widehat w_1(1)=e^{1.68}\cdot0.184\approx0.83$ is \code{U14}/\code{UNIF2}.
3: \code{U11.md} \S1, \code{LOOP-RUNBOOK.md}:19257--19263.
4: \code{U12.md}; \code{LOOP-RUNBOOK.md}:19295--19298.
5: \code{U10.md} \S3 (factor $2.0975$, window $[0.690,1.448]$); acid test at $\beta=0$: \code{LOOP-RUNBOOK.md}:19301--19302.
6: \code{U13.md} \S3 (Lorentzian witness, $f''$ sign change at $2v/\sqrt3$); \code{LOOP-RUNBOOK.md}:19136--19144.
7: \code{MON3.md} \S3(iv) ($\delta_A\cdot a=0.5686159$ versus $a\Lambda(1/48)=0.5700$, ratio $0.9976$; other $\beta$: $1.208/7.394\times10^3/5.188\times10^{10}$); free field $2\pi/8=0.7854$, \code{MON3.md}:185.
8: \code{CB4.md} \S1 row 4 and \code{LOOP-RUNBOOK.md}:16589, :16622; $\Phi(0)=\beta=2N^2u$ is \code{C05-AUDITCBAR.md}:18,52.
9: \code{U10.md} \S1 (J5, threefold redundancy); Fekete side \code{RS5.md} J1, \code{P4-OTHER.md}:145.
10: \code{U10.md} \S\S2--3; \code{U15.md} row 10.
11: \code{U15.md} row 11; $\operatorname{ess\,inf}G''<0$ with $K_S(2.5)=-119$: \code{GAP2.md}:72.
12: \code{U11.md} \S2; witness \code{LOOP-RUNBOOK.md}:19241--19252.
13: \code{U04.md} \S3 (\code{LOOP-RUNBOOK.md}:19180--19194).
14: \code{U13.md} \S1--2 ($2.5/3.777\times10^{-3}=661.9$, i.e.\ $662\times$, \code{U13.md}:56); \code{LOOP-RUNBOOK.md}:19120--19133.
15: \code{LOOP-RUNBOOK.md}:18914--18916 (\S159.1); \code{Q03-TARGET.md}:127; 2D $m=3g^2/8$ and $\nu_{\frac14{\rm tr}F^2}\equiv0$: \code{N02-FOURMODEL.md}, \code{LOOP-RUNBOOK.md}:19088.
16: \code{M05-RATIO.md}:97--100; circularity \code{A04-RATIO2.md}:164 (\code{LOOP-RUNBOOK.md}:18894).
17: \code{Q6-H1-STATUS.md}:208 (thin atom); \code{R6-ATOM-ADVERSARIAL.md}:254 (gapped, atomless); \code{LOOP-RUNBOOK.md}:19048--19053.
18: \code{MON3.md} \S1; \code{LOOP-RUNBOOK.md}:19391 ($dM/d\ln a=0$, $\gamma_m=0$).
19: \code{U02.md}, \code{MON1.md} (\code{LOOP-RUNBOOK.md}:18822--18824, :19300).
20: \code{U16.md} \S3.2, table at :157--164 and :172 (\code{LOOP-RUNBOOK.md}:19351--19355).
21: \code{U09.md} \S3 (\code{LOOP-RUNBOOK.md}:19460--19464, :19490).
22: \code{U09.md} \S3 (\code{LOOP-RUNBOOK.md}:19467--19471).}

Two notes on the table, both honest bookkeeping rather than refinement of the content.
First, the three group headings are not a partition of the \emph{methods} but of the
\emph{outputs}: route 6 proves a true theorem and its hypothesis is false, and so does
route 2, which is why the direction of the bound, not the truth of the theorem, is the
discriminating property. Second, route 14 is graded \LN{} here on the strength of the
interpolation step (\code{LOOP-RUNBOOK.md}:19127--19133), while the summary ledger of that
sweep grades the leg \LC{} with the circularity as its cause
(\code{LOOP-RUNBOOK.md}:19012). The two gradings are both recorded; this paper does not
silently merge them.

\noindent\textbf{The inequalities behind the table, cited.} The correlation inequalities behind rows 2
and 10 are Griffiths \cite{griffiths}, Kelly--Sherman \cite{ksherman}, FKG \cite{fkg} and Ginibre
\cite{ginibre}; log-concavity enters through Pr\'ekopa \cite{prekopa} and Brascamp--Lieb \cite{bl};
row 12 is Temple's bound \cite{temple} together with the Weyl and Lieb--Thirring inequalities; the
``output is an upper bound'' pattern of rows 3--6 is what the transport--entropy criterion
\cite{vrs} and the convolution and entropy-power inequalities \cite{royen,gsv} also produce, and the
negative-dependence theory of \cite{bbl} is what makes a naive FKG substitute unavailable for a
lattice gauge measure; and the rates of rows 5, 7 and 9 exist as stated only inside the
strong-coupling window, which is the log-Sobolev and mixing literature of
\cite{js,sinclair,feller,bbs}. Every one of these is a true theorem. The column that matters is
\emph{what it gives}, not whether it is true.

\subsection{The three layers of the obstruction}\label{sec:survey-three}

The twenty failures are not twenty independent accidents. They decompose into three
layers, each of which is independently sufficient to stop the route, and all three point
the same way.

\paragraph{1. The instrument layer.} The weak-$*$ instruments --- moments, $\varphi^U$,
$a_n$, and every other quantity the campaign can actually compute --- cannot resolve an
atom. The decisive count is $0/14$: of fourteen candidate instruments tested against this
question, none works, and the failure is located in the class of instruments (finite data,
weak-$*$ continuous), not in infinite dimensionality. The mechanism is exact: every
computable quantity is a weak-$*$ continuous functional of $\nu$, whereas $w_1$ is the
mass of a point. A continuous functional cannot see a point mass, so no sequence of
computable quantities converges to $w_1$ in a way that could certify a floor. The one
non-weak-$*$ instrument available is $H1$, and $H1\not\Rightarrow m>0$ (row 17).

\paragraph{2. The continuity layer.} The map $\nu\mapsto\nu(\{x\})$ is only upper
semicontinuous, and a floor requires lower semicontinuity. Concretely, $w_1$ is
\emph{discontinuous} in the K\"allen--Lehmann weak topology: there is a family $\nu_\varepsilon$
with $w_1(\nu_\varepsilon)\to\varepsilon\to0$ across a jump, so that
\[
\inf\;=\;0 \qquad\text{while}\qquad w_1>0\ \text{pointwise.}
\]
Upper semicontinuity permits exactly this downward jump. The quantitative form is the same
statement: for any weak-$*$ continuous $g$ the error obeys $\sup|E-g|\ge f_1^2/2$, so a
continuous instrument cannot even approximate the atom to better than half its mass.
Consequently the target of \S\ref{sec:survey-target} \emph{may be false}, and its being
false is not a defect of the route but a property of the object. The sentence is about $\nu$ as an
unrestricted measure on $[0,\infty)$; read on the atom route of \S\ref{sec:reduction-equiv} the
same statement is the target itself, and that is the entire content of the qualifier printed there.

\paragraph{3. The implication layer.} Even granting the atom, atom-ness does not imply a
gap. The two solvable models separate the implications in opposite directions: 2D
Yang--Mills has a gap and no atom ($m=3g^2/8>0$ while $\nu_{\frac14{\rm tr}F^2}\equiv0$),
and the M6 model has an atom and no gap ($w_1=0.5$ while $m=0$). The stronger statement is
available in 2D and is the cleanest single line in the survey: there $m(L)=\sigma L\to\infty$
while $w_1=0$, so the gap \emph{diverges} while the atom fraction is identically zero.
$w_1$ and $m$ are therefore completely decoupled --- not merely non-proportional.

\begin{v3box}
\textbf{The three layers are independent, and the direction is common.} Layer 1 says the
instruments cannot see the object; layer 2 says the object's own topology forbids the
floor the target asks for; layer 3 says the floored object would not imply the goal even
if floored. No one of the three presupposes another, so removing one obstruction would
leave the other two intact. And the direction of the twenty routes is uniform:
\textbf{upper bounds are in ample supply, every lower bound is circular or
counterexampled.}
\src{\code{U15.md} (the $0/14$ count and the instrument-class localisation);
\code{U08.md} ($\nu_\varepsilon$); \code{R3} via \code{U11.md} \S3 ($\sup|E-g|\ge f_1^2/2$);
\code{U01.md} (usc versus lsc); \code{U06.md} and \code{LOOP-RUNBOOK.md}:19520--19537
(2D, $m(L)=\sigma L\to\infty$ with $w_1=0$); \code{Q03-TARGET.md}:127 and
\code{LOOP-RUNBOOK.md}:19085--19090 (M6); summary \code{LOOP-RUNBOOK.md}:19164--19170,
:19316--19325, :19569--19578.}
\end{v3box}

\subsection{What survives}\label{sec:survey-survives}

Two items survive the sweep, and one of them is a bug.

\begin{itemize}
\item \textbf{A block-Loewner inequality, \LT.} On the dual side the Griffiths/Bochner
  positivity gives $c_r(\beta)\ge0$ and $\beta$-non-decreasing (tensor-product
  multiplicities are non-negative; verified numerically on SU(3) at $\beta=2.5$ and
  $5.7$), and hence $K_\beta\succeq K_{\beta'}$ for $\beta\ge\beta'$, for all $N$. This
  repairs the previously unproved step of D05. It does not reach $w_1$: it is blocked by
  the $D$-factor and it does not transport eigenvectors.
\item \textbf{A sampler bug, \LE.} The same sweep found that the QR sampler fabricates
  sign changes. Four corpus conventions therefore require revision, and the earlier
  numerical determinations of ``$c_r$ changes sign'' are void. This is a by-product of a
  negative survey, and it is the only place in the paper where a negative result removed a
  positive claim rather than failing to establish one.
\end{itemize}

The last observation of the survey closes the FKG family as well, and closes it at $H1$.
The only monotonicity $w_1$ admits is on the three spectral points $(s_1,s_2,\bar E)$, in
the right direction:
\[
w_1\;\ge\;\frac{s_2-\bar E}{s_2-s_1},
\]
and $w_1$ is not monotone in $\beta$ (a $\tfrac14\sin^2\theta$ hump) and is constant under
$O\mapsto\lambda O$ ($\gamma=0$). The threshold form of that inequality is
\[
w_1\ge w_0 \iff \frac{\delta_1}{\varepsilon}\le 1-w_0,
\]
which is precisely the isolation of the second support point, i.e.\ $H1$. So the FKG route
does not deliver a floor either: it delivers $H1$, and $H1$ does not imply $m>0$ (row 17).
Combined with the instrument layer, the conclusion is that positivity in this family reaches
the second support point and stops there.

\src{\code{U09.md} \S3 and \code{LOOP-RUNBOOK.md}:19460--19464, :19467--19471, :19490--19498
(block-Loewner, sampler bug, and the reduction $w_1\ge w_0\iff\delta_1/\varepsilon\le1-w_0$ =
second support point $=H1$); $H1\not\Rightarrow m>0$: \code{Q6-H1-STATUS.md}:208,
\code{LOOP-RUNBOOK.md}:19048--19053.}

% =====================================================================
%  sec05-ledger.tex --- Section 5: the bookkeeping ledger (38 entries)
%  Merged edition v3.  All-English body.  Every number carries file:line.
% =====================================================================

\section{The bookkeeping ledger}\label{sec:ledger}

\subsection{Why a ledger is in a mathematics paper}\label{sec:ledger-why}

A table of one's own errors earns its place in a paper only if it does work. This one does: it is the
mechanism by which the results of this paper were corrected, not an appendix to them. Three of the five
examples below moved a numeral that appears in the body; the other two stopped a sentence before it
reached a page. Without the ledger the same conclusion becomes available again on the next pass, because
nothing in the mathematics records that it was already tried and failed.

\begin{claimbox}{cN}
\claimhead{Claim}{.} Each of the five entries below is load-bearing: deleting its row either changes a
numeral in this paper or restores a sentence that was withdrawn from a draft. The ledger is kept for that
reason alone, and not as a record of diligence.
\src{\code{LOOP-RUNBOOK.md} lines 9986, 12111, 16768, 17735, 19672;
\code{ADJ6.md} lines 1, 11, 135, 138; \code{Q13-BETA6.md} line 59;
\code{position-current.html} lines 412--414.}
\end{claimbox}

\noindent\emph{Path convention.} File names are given without their directories throughout this section:
the position file lives in \code{docs/}, the species files in \code{research/v3/species/}, and the
chapter files in \code{docs/paper-v2/}. Line numbers are line numbers in those files, and two
abbreviations are used for the two files that recur: \code{RB:17735} is line 17735 of
\code{LOOP-RUNBOOK.md}, \code{P:378} is line 378 of the position file. Five further short labels are
carried over verbatim from the corpus --- \code{V6}, \code{p8}, \code{F07}, \code{ADJ6} and
\code{Q13} --- because the entries that use them were quoted from the chronicle rather than
re-typed; this paper does not resolve them to paths, and a reader who wants them should grep the
label in \code{LOOP-RUNBOOK.md}. Naming the gap is preferable to inventing a path for it.

\begin{enumerate}[label=\textcircled{\small\arabic*},leftmargin=2.3em,itemsep=4pt,topsep=4pt]
\item \textbf{Four sequences, one label.} The sequences $a_n$, $B_n$, $m_n$, $t_n$ are quoted with two
labels for one object and with one label for two different objects. Whether $\lambda_T$ is bounded
through $t_n$ or through $a_n$ decides whether the shortfall against $48$ is $4.75\%$ or $13.3\%$ --- one
physical statement, two numbers. The row settles it: $B$ bounds both sequences, the shortfall is
$B/48=1.0475$, and $13.3\%$ is the branch computed with a dominated constant.
(\code{ADJ6.md:1,11,135,138}; \code{LC3.md:21}; \code{RB:17735}.) This is row~12 of
\S\ref{sec:ledger-thirtyeight}.

\item \textbf{$W_{\mathrm{full}}$ versus $W_{\mathrm{geo}}$: one name, two objects.} The two were used
as the same object for several rounds. Pr\'ekopa's inequality \cite{prekopa} --- and its
Brascamp--Lieb form \cite{bl} --- gives $W_{\mathrm{geo}}\ge0$; it says
nothing about $W_{\mathrm{full}}$. A row that merged them transfers a theorem to an object the theorem
does not concern. (\code{RB:9986}.) This is row~10.

\item \textbf{A first moment that is a second moment.} $\bar R=M_1/M_0$ in \code{p8} was reported as a
first time moment; it is $\langle m^2\rangle$. The power is off by one, and the bridge that consumed the
number was consequently reported stronger than it is. The row is the reason that bridge carries no \LT{}
anywhere in this paper. (\code{RB:19672}.) This is row~19.

\item \textbf{A paper number on the wrong author.} Numbers $107/110$ were attributed to Balaban; they
are Federbush's, and a three-source search for $113/115/117$ returned $0$ hits. The wrong attribution had
already entered a draft, which is the only reason it is a ledger row rather than a correction made in
passing. (\code{RB:12111}.) This is row~23.

\item \textbf{A number detached from its coupling.} The value $0.6018$ is the pooled result of the
$\beta=6$ run ($\chi^2/\mathrm{dof}=1.22$, eight groups consistent). Quoted without $\beta=6$ it becomes
a number at $\beta=2.5$, which is the coupling of the measurement in this paper.
(\code{Q13-BETA6.md:59}.) This is row~9.
\end{enumerate}

\begin{v3box}
\textbf{Three further entries reached a draft and were withdrawn.} These are not hypothetical: each was
written as a conclusion, not as a suspicion.
\begin{itemize}[leftmargin=1.5em,itemsep=3pt,topsep=3pt]
\item $\Delta_0>0$ was written as the name of the residual gap and withdrawn. Verbatim: ``The earlier
naming of that gap as $\Delta_0>0$ is WITHDRAWN (n12): $\Delta_0$ sits on $G''$, not on the measure
$\nu$.'' (\code{paper-v6/main.tex:17}.) This is the object-mismatch class, row~15, caught at draft
level: the name sat on a carrier that is not the carrier of the statement.
\item $0.34/0.46$ was written into a draft as the $\mathrm{SU}(2)/\mathrm{SU}(3)$ values
(\code{ch05-theorems.tex:208--209}). The exact $\mathrm{SU}(2)$ Haar-slice value is
$0.20512$; $0.34/0.46$ are the magnitudes for $N=3/N=4$ ($0.342/0.461$). The numbers survive, the label
does not (\code{O09-PHICURV.md:104}). This entry has \emph{no} itemized identity in
\S\ref{sec:ledger-thirtyeight}: its only trace in the corpus is the correction itself, and it is
therefore not counted among the twenty-nine rows.
\item $W_{\mathrm{full}}$ and $W_{\mathrm{geo}}$ were used as one object in a draft, i.e.\
example~\textcircled{\small2} at draft level. Row and draft error are one entry seen twice, and they are
counted once. (\code{RB:9986}.)
\end{itemize}
\end{v3box}

\noindent\textbf{Why this is not a repository log.} A reader who re-uses $w_1(2.5)$, or the pooled
$0.6018$, or the window $\beta_A<3.777\times10^{-3}$, needs to know which row sits under the number. In a
body of results whose only positive content is a measurement, the audit trail is part of the mathematics:
it is what makes the number quotable at all, and it is what a referee is entitled to check.

% =====================================================================
\subsection{The thirty-eight}\label{sec:ledger-thirtyeight}

Two counting rules, stated before the count, because a different rule set yields a different total and
the total is the thing this paper has to defend. \emph{One entry} is one minimal unit of debt such that
copying it literally produces a wrong conclusion. A polysemous symbol counts once per symbol, with the
number of senses in brackets; a clash of two numerals or of two sequences counts once per confused pair.
The rules do not answer ``have all consequences of the debt been settled'' --- each row carries its own
consequence and its own source.

{\small
\setlength{\tabcolsep}{4pt}
\begin{longtable}{@{}>{\raggedright\arraybackslash}p{0.055\textwidth}%
>{\raggedright\arraybackslash}p{0.163\textwidth}%
>{\raggedright\arraybackslash}p{0.505\textwidth}%
>{\raggedright\arraybackslash}p{0.207\textwidth}@{}}
\caption{The thirty-eight entries: class, the one sentence that states the debt, and what the entry
moves. Rows 30--38 are a class-level residual with no itemized identity
(\S\ref{sec:ledger-thirtyeight}, Gap~1).}\label{tab:ledger}\\
\toprule
\textbf{\#} & \textbf{Class} & \textbf{One sentence} & \textbf{Grade impact} \\
\midrule
\endfirsthead
\toprule
\textbf{\#} & \textbf{Class} & \textbf{One sentence} & \textbf{Grade impact} \\
\midrule
\endhead
\midrule
\multicolumn{4}{r@{}}{\footnotesize continued on the next page}\\
\endfoot
\bottomrule
\endlastfoot
% ------------------------- deepest -----------------------------------
\textbf{1}\newline{\scriptsize item 29} & Symbol definition, non-unique (deepest) &
$u$ is defined twice, $u:=c_f$ and $u:=c_f/N$, and the two differ by $N$; every formula containing $u$
therefore inherits the ambiguity. (\code{RB:16768}.) &
No local grade moves; it sits under the window, $\bar C$, $\bar u$, $r_0$ and $\delta$. \\
\midrule
% ------------------------- polysemy 6 --------------------------------
\textbf{2} & Symbol polysemy (1/6) & $\bar C$ carries three senses (sum-level constant, condensate,
coefficient); one of them is the identity $\bar C=1$. & Withdraws $\bar C=1$; the $n=3$ witness stands
(\code{V6:14}). \\

\textbf{3} & Symbol polysemy (2/6) & $\delta$ carries six senses: step size, mass lower bound,
certificate margin among them. & No grade moves; makes the comparison of row~21 possible at all. \\

\textbf{4} & Symbol polysemy (3/6) & $\Omega$ carries four senses. & None recorded; the
$\Omega$/Kato route is already circular (\S\ref{sec:ledger-dead}). \\

\textbf{5} & Symbol polysemy (4/6) & $G$ is written both for a propagator and for the potential that
$G''$ differentiates. & None recorded; every $G''$ statement must name its carrier. \\

\textbf{6} & Symbol polysemy (5/6) & $E[H]$ carries two senses: the energy of a configuration and an
expectation under the measure. & None recorded; both readings are in play in the M6 test. \\

\textbf{7} & Symbol polysemy (6/6) & $\Delta_{\mathrm{dyn}}$ carries two definitions, one of them
$\Delta_{\mathrm{dyn}}=-W_{\mathrm{full}}$ on 2D Yang--Mills. & Blocks the four-model test until one
definition is fixed. \\
\midrule
% ------------------------- numerical clash 4 -------------------------
\textbf{8} & Numerical name clash (1/4) & $1/48$ and $1.766/48$ are both called ``the threshold'';
the first is the P3 residue constant, the second the shifted one. & The $\beta_A$ window inherits the
factor $1.766$. \\

\textbf{9} & Numerical name clash (2/4) & $0.6018$ is the pooled $\beta=6$ value
($\chi^2/\mathrm{dof}=1.22$); detached from $\beta=6$ it reads as a $\beta=2.5$ number. &
Quotable only at $\beta=6$ (\code{Q13-BETA6.md:59}). \\

\textbf{10}\newline{\scriptsize item 5} & Numerical name clash (3/4) & $W_{\mathrm{full}}$ and
$W_{\mathrm{geo}}$ are two objects under one name, used as one object. & Pr\'ekopa's
$W_{\mathrm{geo}}\ge0$ does not transfer to $W_{\mathrm{full}}$ (\code{RB:9986}). \\

\textbf{11} & Numerical name clash (4/4) & $\delta_A\cdot a$ and $a\Lambda(1/48)$ differ by $0.24\%$,
and the agreement is read as an identity. & No grade moves; the $0.24\%$ is not an identity test. \\
\midrule
% ------------------------- sequence clash 2 --------------------------
\textbf{12}\newline{\scriptsize item 30} & Sequence name clash (1/2) & $a_n$, $B_n$, $m_n$, $t_n$ are four
sequences under two labels; $t_n$ versus $a_n$ decides $4.75\%$ versus $13.3\%$. & $13.3\%$ is
withdrawn; $4.75\%$ survives (\code{ADJ6.md:11,135,138}; \code{RB:17735}). \\

\textbf{13} & Sequence name clash (2/2) & The old $T^4_6$ convention counts $7218758$ at $n=6$, short by
$10$; the clean value under $nx\ge8$ is $7218768$. & Any count built on the old convention is short by
$10$ (\code{RB:18570}). \\
\midrule
% ------------------------- scope error 5 -----------------------------
\textbf{14} & Scope error (1/5) & The defect ratio $D/S\approx2\%$ at $\beta=2.5$ is quoted without its
denominator. & No grade moves; the $2\%$ needs the word ``defect'' attached (\code{P:414}). \\

\textbf{15} & Scope error (2/5) & $\Phi$ is attached to the wrong carrier: $\Delta_0$ sits on $G''$,
not on the measure $\nu$. & The name $\Delta_0>0$ is withdrawn (\code{V6:17}). \\

\textbf{16} & Scope error (3/5) & The input $|w|^{1/n}=c_f$ at \code{F07:82} rests on $|w_p|\le u$,
which is wrong by the factor $N$. & Removes one input to the $\bar C$ route (\code{V6:14}). \\

\textbf{17} & Scope error (4/5) & $a=0$ is an activation point of the recursion, not a criterion; using
it as one over-reports the window by $4.84\times$. & Window corrected to
$\beta_A<3.777\times10^{-3}$ from $0.018394$ (\code{main.tex:162--163}). \\

\textbf{18} & Scope error (5/5) & $Q_3$ is misread: the vertex count of the cube is quoted as a slice
count of the decomposition. & No grade moves; the count must be recomputed per object. \\
\midrule
% ------------------------- power error 3 -----------------------------
\textbf{19}\newline{\scriptsize item 34} & Power error (1/3) & $\bar R=M_1/M_0$ in \code{p8} is reported as a
first time moment; it is $\langle m^2\rangle$, one power off. & The bridge is over-stated and carries no
\LT{} (\code{RB:19672}). \\

\textbf{20} & Power error (2/3) & $\beta$ and $\beta^2$ appear in one comparison under one symbol. &
Any inequality mixing the two is void. \\

\textbf{21} & Power error (3/3) & The contradiction $.577>.551$ in \code{U14} compares two different
$\delta$'s. & No grade moves; it is an artefact of row~3. \\
\midrule
% ------------------------- citation / nomenclature 4 -----------------
\textbf{22} & Citation, nomenclature (1/4) & A citation is rewritten in paraphrase and then quoted as
if verbatim. & The cited claim's grade is not inherited. \\

\textbf{23}\newline{\scriptsize item 14} & Citation, nomenclature (2/4) & $107/110$ are attributed to Balaban;
they are Federbush's, and $113/115/117$ return $0$ hits. & Attribution corrected; the Balaban chain is
not available (\code{RB:12111}). \\

\textbf{24} & Citation, nomenclature (3/4) & ``RG stable modulus'' occurs $4$ times and is defined $0$
times. & The RG-modulus route has no stated object. \\

\textbf{25} & Citation, nomenclature (4/4) & $w_1$ is defined with no $t$ axis, so $w_1(2.5)$ is a
column at fixed $t$. & $w_1(2.5)$ stays \LE; no $t$-limit is available. \\
\midrule
% ------------------------- new this round 4 --------------------------
\textbf{26}\newline{\scriptsize item 35} & New this round: denominator outside its window & The denominator
$r_0=1/32$ of the $10.94\times$ is negative at $\beta=2.5$ ($r_0$ is a floor only for
$\beta\le0.147152$; $\bar u(2.5)=16.99$ gives $r_0<0$). & The ``floor'' reading is withdrawn; the
$10.94\times$ is a dimensionless ratio, not a proved floor (\code{RB:19728--19730}; \code{RB:19738}). \\

\textbf{27}\newline{\scriptsize item 36} & New this round: over-generalized sign statement & ``\code{APE2}
and \code{APE4} show $C(2)<0$'' is too broad: only $L=6$ \code{APE4} at $-0.00740$ and \code{APE1} at
$-0.00121$ are negative; there is no negative \code{APE2} row. & The sign statement is narrowed to
$L=6$, \code{APE1}/\code{APE4} (\code{RB:19775--19779}; \code{RUN1:48,51}). \\

\textbf{28}\newline{\scriptsize item 37} & New this round: grade demotion & Step (iii) of \code{SG1-1} holds
only under a fixed-$N$, continuity reading; read as computability, the three-line proof fails. &
\code{SG1-1} is \LC{} and not \LT, after being labelled \LT{} for $60+$ rounds
(\code{RB:19797--19812}). \\

\textbf{29}\newline{\scriptsize item 38} & New this round: quantifier scope & The target uniqueness of
\code{N09-1} holds only at fixed $O$: $w_1$ depends on $O$, $0.343$ versus $0.017$, a factor $20.2$. &
\code{N09-1} is restricted to fixed $O$; the uniqueness is not $O$-free
(\code{RB:19849--19879}; \code{P:378}). \\
\midrule
% ------------------------- residual 9 --------------------------------
\textbf{30} & Class-level residual (1/9) & No itemized identity in any source: the itemized list is
truncated by an ellipsis at \code{position-current.html:414}. & Counted, not quotable. \\

\textbf{31} & Class-level residual (2/9) & No itemized identity in any source: the itemized list is
truncated by an ellipsis at \code{position-current.html:414}. & Counted, not quotable. \\

\textbf{32} & Class-level residual (3/9) & No itemized identity in any source: the itemized list is
truncated by an ellipsis at \code{position-current.html:414}. & Counted, not quotable. \\

\textbf{33} & Class-level residual (4/9) & No itemized identity in any source: the itemized list is
truncated by an ellipsis at \code{position-current.html:414}. & Counted, not quotable. \\

\textbf{34} & Class-level residual (5/9) & No itemized identity in any source: the itemized list is
truncated by an ellipsis at \code{position-current.html:414}. & Counted, not quotable. \\

\textbf{35} & Class-level residual (6/9) & No itemized identity in any source: the itemized list is
truncated by an ellipsis at \code{position-current.html:414}. & Counted, not quotable. \\

\textbf{36} & Class-level residual (7/9) & No itemized identity in any source: the itemized list is
truncated by an ellipsis at \code{position-current.html:414}. & Counted, not quotable. \\

\textbf{37} & Class-level residual (8/9) & No itemized identity in any source: the itemized list is
truncated by an ellipsis at \code{position-current.html:414}. & Counted, not quotable. \\

\textbf{38} & Class-level residual (9/9) & No itemized identity in any source: the itemized list is
truncated by an ellipsis at \code{position-current.html:414}. & Counted, not quotable. \\
\end{longtable}
}

\noindent\emph{Reading the table.} The classes add up as stated in the corpus: $1$ deepest $+$ $6$
polysemy $+$ $4$ numerical $+$ $2$ sequence $+$ $5$ scope $+$ $3$ power $+$ $4$ citation $=25$ named
entries, plus the $4$ added this round $=29$, plus the $9$ class-level residual $=38$. The three
examples most often re-quoted ($w_1(2.5)$, $0.6018$, $\beta_A$) sit under rows 25, 9 and 17 respectively.

\begin{warnbox}
\textbf{Gap 1: twenty-nine of thirty-eight.} Of the $38$ entries, $29$ carry an itemized identity ---
$25$ named one by one and $4$ added this round (rows 26--29). The remaining $9$ (rows 30--38) exist as a
class-level residual only: the corpus states the total, and its itemized list is cut off by an ellipsis
at \code{position-current.html:414}, where the visible list contains $24$ entries beside item $29$.
$29+9=38$; the arithmetic of the total is therefore checked, and the identity of nine entries is not.

\textbf{Gap 2: the label sequence has holes.} Labels of the form ``item $N$'' in
\code{LOOP-RUNBOOK.md} cover only $30$ of the numbers $1$--$34$: the labels $1$,
$2$, $4$ and $26$ never occur. Inside the table above this is visible as absence: rows 1, 10, 12, 19 and
23 carry a recovered ledger number in the \#{} column, the other named rows do not, because the corpus
does not give one.

\textbf{The nine are not manufactured.} A ledger that invented nine entries in order to reach a round
number would commit precisely the error it exists to prevent. The residual is reported as a residual, and
any future reader who wants those nine must find them in the chronicle, not in this table.
\end{warnbox}

\subsubsection*{Verification of the two counts (search recorded, not asserted)}

Per the provenance rule of this paper, the searches behind the two gaps are recorded with tool state,
query and hit count, so that ``0 hits'' is never confused with ``the tool failed''.

\begin{table}[htbp]
\centering\small
\setlength{\tabcolsep}{4pt}
\begin{tabularx}{\textwidth}{@{}>{\raggedright\arraybackslash}p{0.40\textwidth}%
>{\raggedright\arraybackslash}p{0.08\textwidth}%
>{\raggedright\arraybackslash}X%
>{\raggedright\arraybackslash}p{0.13\textwidth}@{}}
\toprule
\textbf{Query} & \textbf{Tool} & \textbf{Result} & \textbf{Hits} \\
\midrule
regex \code{\textbackslash u7b2c(\textbackslash d\{1,2\})\textbackslash u5904} over
\code{LOOP-RUNBOOK.md} & grep & valid response; $34$ distinct numbers, $3$, $5$--$25$, $27$--$38$ &
$51$ lines \\
the same, restricted to $1$--$34$ & grep & valid response; $30$ present, $1$, $2$, $4$, $26$ absent &
$0$ labels \\
\code{38} beside the counter token, repository-wide & grep & valid response;
\code{position-current.html:412,420} & $2$ lines in that file \\
\bottomrule
\end{tabularx}
\caption{The two count checks, recorded with tool state, query and hit count, so that ``$0$ hits'' is
never confused with ``the tool failed''.}\label{tab:ledger-verify}
\end{table}

\noindent Two readings follow. First, $30\le34$ is consistent: the label sequence is a subset of the
ledger, so the missing four neither raise nor lower the total. Second, the four missing numbers are
\emph{not} claimed to be empty entries --- whether entries with those numbers exist is something this
paper does not assert, and the count does not depend on it.

% =====================================================================
\subsection{The dead-route ledger}\label{sec:ledger-dead}

The second ledger is given as classes and counts only: the individual routes are surveyed elsewhere in
this paper, and itemizing them a second time would duplicate that survey without producing a number. All
counts below are from \code{position-current.html:409}, where the same table appears as ``the
dead-route ledger, $35+$ entries, one number each''.

\begin{table}[htbp]
\centering\small
\setlength{\tabcolsep}{4pt}
\begin{tabularx}{\textwidth}{@{}>{\raggedright\arraybackslash}p{0.36\textwidth}%
>{\raggedright\arraybackslash}p{0.07\textwidth}%
>{\raggedright\arraybackslash}X@{}}
\toprule
\textbf{Class} & \textbf{Count} & \textbf{Comment} \\
\midrule
Circularity, or a direction reversed ($\Rightarrow$ read as $\Longleftrightarrow$) & $9$ &
$H1\Rightarrow m>0$, $I2\Rightarrow H1$, $H1\Rightarrow I2$, $I2\Rightarrow m>0$ bare,
$I2\Longleftrightarrow\rho^2>1$, the \code{L1} chain, the \code{L2} reduction, $\Omega$/Kato/$\det>0$,
transport \\
Structural counterexample & $5$ & window enlargement self-refuting, $(\star)$ uniform, the sixth
species, the ratio lemma ($\max b(4)=62$), the layered decomposition \\
Numerically or by scale unreachable & $5$ & \code{GEVP} ($2.3\times10^7$ configurations, $310$\,h), the
$u$-series at $2.5$, $\hat K=8$ from six directions, the \code{P3} loop-order grading, the $19$-item
transport ledger \\
Sufficiency misjudged, or tautology & $4$ & $G''$ sufficient, $f_1^2$ necessary,
$W^{1/2}W^{-1/2}$ similarity-invariant, \code{SOS} with $\lambda_{\min}=-g|B|-c$ \\
Proxy, or substitution & $2$ & $\Delta_{\mathrm{dyn}}$ as proxy, the three machines \\
Taxonomy, scope, consistency & $3$ & the five classes, $1/48$ preserved as it stands, the twenty-route
consistency check \\
\midrule
\textbf{Itemized-checkable subtotal} & $\mathbf{28}$ & $9+5+5+4+2+3$ \\
Class-level residual (claimed $35+$, not itemized) & $\ge7$ & computed as claimed $35$ minus itemized
$28$ \\
\midrule
\textbf{Total} & $\mathbf{\ge35}$ & consistent with the ``$35+$'' claim at \code{P:409} \\
\bottomrule
\end{tabularx}
\caption{The dead-route ledger as classes and counts. Itemized-checkable $28$, claimed $35+$, and the
difference stored as a class-level count rather than as invented rows.}\label{tab:deadroutes}
\end{table}

\noindent The same shape as \S\ref{sec:ledger-thirtyeight}: itemized-checkable $28$, claimed $35+$, and
the difference $\ge7$ stored as a class-level count only, not as invented rows. The directional reading of
the same table is at \code{position-current.html:397}: nine routes give upper bounds only, nine are
circular, two have explicit counterexamples, two produced something usable. That is the campaign's failure
written as a shape --- upper bounds are complete, and every lower bound terminates in the target.

\begin{qbox}
\textbf{What the two ledgers do not claim.} Neither ledger claims that the consequences of a debt are all
settled; each row answers one question only, namely whether the debt exists and what it moves. The
dead-route counts are not adjusted for item~37: the sixth species still dies of \code{SG1-1}, but
\code{SG1-1} is \LC{} rather than \LT{} this round, so the kill is conditional and the row 28 qualifier
travels with it. No total in this section is a theorem, and none of them is offered as progress.
\end{qbox}

% =====================================================================
%  PART VI -- THE VERDICT, THE DISCLOSURE, AND THE APPENDIX
%  (merged edition v3, section 6)
% =====================================================================
\clearpage
\part{The verdict, the disclosure, and the appendix}

\section{The verdict}
\label{sec:verdict}

\subsection{Clay: \texorpdfstring{$0\%$}{0 percent}}
\label{sec:verdict-clay}

This subsection states the figure first and the reasons second, because the figure is the
deliverable and the reasons are only its justification.

\begin{claimbox}{cN}
\textbf{Clay: $0\%$.} The four-dimensional pure $SU(N)$ Yang--Mills mass gap is not proved in
this paper, is not reduced to a finite computation in this paper, and is not placed within one
lemma of any tool developed here. What is proved is a reduction, a set of supporting theorems, one
measurement at the physical coupling, and a systematic negative survey. The percentage above is a
judgement about distance to a complete unconditional proof, not a measure of effort.
\end{claimbox}

The honest structure of the failure is not ``we did not find the argument''. It is that
\textbf{each of the three legs of the reduction closes on its own premise}: the reduction is exact,
and every attempt to discharge a leg returns that leg. We take the three legs in turn.

\subsubsection*{(i) Pointwise positivity at the physical coupling}

The leg asks for $w_1(2.5) > 0$ for the physical operator. \textbf{Numerically, it holds.} On
$L_t = 12$, with the \code{njit} heat bath at $\beta = 2.5$, on $L_s = 6$ and $L_s = 8$, the
measurement gives
\[
  w_1(2.5) \;=\; 0.3427 \,\pm\, 0.0119 \qquad (28.8\,\sigma \text{ from zero}),
\]
so the sign is not in doubt at this lattice spacing and this volume. But the statement is not a
statement about the Hamiltonian: on the \emph{same} Markov chain and the same $12000$
configurations, a second operator gives $0.0168 \pm 0.0142$. \textbf{The two differ by a factor
$20.2$.} The one-line counterexample $w_1 = \cos^2\theta$ at fixed spectrum shows why: $w_1$ is a
property of the pair $(H, O)$, not of $H$. Hence the leg splits in two: its pointwise form is \LE{}
(numerically satisfied), while its uniform form is \LO{} and, until $O$ is named, vacuous as a
statement about the theory. Transporting it to the continuum requires an $O$-independence statement
that does not exist here.

\subsubsection*{(ii) The $L$-monotone step}

The leg asks for a monotone quantity in the volume that survives the infinite-volume limit.
\textbf{No such step exists in the class of quantities considered.} The obstruction is two-sided
realizability: both a $2/(L-1)$ profile and a constant profile are realizable within the same
structural assumptions, so no monotonicity theorem can be extracted from the structure alone; and
there is \textbf{no FKG inequality} for the lattice gauge measure at $\beta = 2.5$ to supply the
comparison that realizability leaves open. The leg is therefore not merely unproved but
unavailable: nothing in the present toolkit produces it. \LN{}

\subsubsection*{(iii) The $\beta$-tail-uniform floor}

The leg asks for a floor on the spectral quantity that is uniform in $\beta$ along the tail.
\textbf{This floor is $\Delta_0$, and $\Delta_0$ has already been judged dead}: the model $M6$ is
insufficient to supply it, and the two-dimensional theory is unnecessary for it, so neither the
solvable-model route nor the dimensional-reduction route reaches it. Renaming the floor does not
create one; the leg and the target are the same object. \LO{}

\begin{v3box}
\textbf{Reading of the three legs together.} (i) holds numerically and depends on $O$ with a factor
$20$; (ii) does not exist as a step, and no FKG is available to create it; (iii) is $\Delta_0$,
already dead. The three failures are of three different kinds --- empirical-but-operator-bound,
structurally-absent, and identically-the-target --- which is why the reduction does not collapse
into a single missing lemma.
\end{v3box}

\src{Leg (i): the measurement of \S\ref{sec:measurement}, $\beta = 2.5$, $L_t = 12$; operator
pair as stated in \S\ref{sec:m-opdep}. Leg (ii): two-sided realizability construction,
\code{research/unconditional/} notes. Leg (iii): $\Delta_0$ and its ledger entry,
\S\ref{sec:ledger}.}

\subsubsection*{The direction of the twenty routes}

The verdict is reinforced, rather than merely asserted, by the survey of \S\ref{sec:survey}: twenty
routes to uniformity in $(L, a)$ were attacked, and they fail in a \emph{uniform direction}. That
uniformity is the substantive finding, not the individual failures.

\begin{center}
\small
\begin{tabularx}{\textwidth}{@{}l X c@{}}
\toprule
\textbf{Class} & \textbf{Routes} & \textbf{Count} \\
\midrule
Upper bound only & compactness, reflection positivity, Berezin--Lieb, the probabilistic reading,
the Krein--Rutman residual, concavity, the RG modulus, two Ursell constants & $9$ \\
Circular & positivity improvement, minimality, operator inequalities, definability, the
strong--weak bridge, and the named obstructions $\Delta_0$, $T4$, $H1$, $M(\beta)$ & $9$ \\
Counterexample & two natural implications, each killed by an explicit machine-checked
counterexample & $2$ \\
\bottomrule
\end{tabularx}
\end{center}

\noindent The direction is one-way: \textbf{the upper side is complete and the lower side is
empty.} Nine routes give only upper bounds on the quantity wanted; nine give lower bounds only by
assuming what is to be proved; two are false as stated. Two routes did produce something usable ---
a block-Loewner inequality that repairs a previously unproved step, and a sampler bug whose
correction invalidates four earlier numerical sign determinations --- but neither is a lower bound.
A survey in which every independent route fails in the same direction is evidence about the shape
of the problem, not a list of unlucky attempts.

\subsubsection*{The only framework that survives, and what it gives}

Exactly one piece of general machinery survives the survey without circularity: the pair
\textbf{Banach--Alaoglu} (for the existence of weak limits of the finite-volume measures) together
with \textbf{Portmanteau} (for the identification of the limit's expectation functionals). Both
are free, and both are free in the direction that does not help.

\begin{itemize}
\item \textbf{Tightness is free.} The lattice measures are tight by construction on the compact
field space, so the compactness class of arguments is excluded as a source of new information: it
supplies existence of a limit, never a gap.
\item \textbf{Upper semicontinuity is free.} Portmanteau then gives the inequality that runs from
finite volume to the limit, and only in that direction. There is no companion inequality running
from the limit back to finite volume, and the gap is exactly the missing companion.
\end{itemize}

\noindent Consequently the free machinery yields precisely the direction that the nine upper-bound
routes already used, and says nothing about the lower side. This is the technical content of the
sentence ``the wall is a direction, not a difficulty''.

\subsection{What would be needed}
\label{sec:verdict-needed}

The reduction of \S\ref{sec:reduction-equiv} identifies the missing object, and it is worth naming
its \emph{type} rather than its content, because every previous attempt in this project named the
content and thereby restated the target.

\begin{qbox}
\textbf{What is missing is one machine that produces lower bounds} --- not another transport of an
upper bound across the same interface. The distinction is not rhetorical: the nine upper-bound
routes of \S\ref{sec:survey} all share a common form (a free inequality applied to a free
compactness or semicontinuity statement), and the nine circular routes all share the complementary
form (a lower bound obtained by assuming a lower bound). A machine of the missing type must produce
a lower bound from hypotheses that do not include one.
\end{qbox}

Every attempt in this project to build that machine terminated at the same point, and it terminated
there \emph{independently} in three places:

\begin{enumerate}
\item \code{N09-1}: the target is unique --- any object strong enough to be transported is already
the gap, so there is no strictly weaker surrogate to transport.
\item \code{u04}: the non-circular version of the argument \emph{is} the Clay problem --- removing
the circular step removes the argument, not the circle.
\item \code{mon1}: the point $L = \infty$ is itself the gap --- the limit object whose positivity
would be transported is the object whose positivity is wanted.
\end{enumerate}

\noindent Three independent lines of work, arriving at three different formulations, identify the
same point. We regard that coincidence as the strongest single piece of structural evidence
produced by the campaign, and it is the reason the verdict is stated as a direction rather than as
a shortage of effort.

\begin{qbox}
\textbf{This subsection is not a restatement of \S\ref{sec:verdict-clay}, and the distinction is
deliberate.} \S\ref{sec:verdict-clay} lists \emph{which} legs are closed and \emph{why} each closure
is of its own kind; this subsection states \emph{what type of object} is absent, and records the
three independent identifications of that object with the target. The first is a status report on
three named legs; the second is a claim about the shape of any future proof. Where the two would
have overlapped --- the naming of $\Delta_0$ as the third leg and as one of the nine circular routes
--- the material is stated once, in \S\ref{sec:verdict-clay}, and referenced here.
\end{qbox}

\subsection{What is nevertheless new}
\label{sec:verdict-new}

The negative verdict does not license self-effacement. Five results below are new, are owned by
this paper, and would survive the disappearance of every other claim in it.

\begin{enumerate}
\item[\textbf{1.}] \textbf{A window correction of a factor $4.87$.} A previously published
criterion window, $0.018394$, is the $a = 0$ activation root and is illegal as a criterion window;
the correct value is $3.777\times 10^{-3}$, and $0.018394/0.003777 = 4.87$ (the companion ratio
$4.84$ belongs to the earlier $0.018269$ of \S\ref{sec:est-window-fix}). The over-report was found
and corrected by this project
against its own published claim. A published over-report that is caught by its own author is a
result, and we do not downgrade it to an erratum.
\item[\textbf{2.}] \textbf{The first numerical value of $w_1$ at $\beta = 2.5$.} The measurement
$w_1 = 0.3427 \pm 0.0119$ at $28.8\sigma$, on $L_s = 6$ and $L_s = 8$ at $L_t = 12$, with the
operator dependence ($0.0168 \pm 0.0142$ for the second operator, factor $20.2$) explicitly
labelled rather than suppressed. The labelling is part of the result: an unlabelled version of this
number would be misleading.
\item[\textbf{3.}] \textbf{The no-go of \code{SG1-1}} (\LC{}). A sixth species of proof is
excluded, at grade \LC{} rather than \LT{} because one step does not survive a computability
reading. The downgrade from \LT{} to \LC{} was reported by the author of that result, not by a
critic.
\item[\textbf{4.}] \textbf{A systematic negative survey of twenty routes} (\S\ref{sec:survey}),
organized by the direction of failure rather than by chronology. Its value is the invariant
direction --- nine upper, nine circular, two counterexamples --- not the count.
\item[\textbf{5.}] \textbf{A ledger of thirty-eight definitional collisions} (\S\ref{sec:ledger}).
This is a methodological contribution: it is the mechanism by which the errors of this paper were
caught, including twice during the writing of this manuscript.
\end{enumerate}

\noindent Item 5 is not bookkeeping decoration. Two of the thirty-eight entries --- corpus item~35,
a negative $r_0$ (row~26 of Table~\ref{tab:ledger}), and corpus item~36, an over-broad formulation
of the asymptotic-perturbation estimate (row~27) --- were found
\emph{while this section was being drafted}, and both were found by an independent review pass
rather than by the author of the affected claim. That is the mechanism, and it is the subject of
the next subsection.

\subsection{AI contribution statement}
\label{sec:verdict-ai}

This paper is a multi-agent artifact, and the statement below is written so that a reader can tell
what was done by which kind of agent.

\begin{claimbox}{cC}
\textbf{The great majority of the derivations, the numerical work, and the audits in this paper were
carried out by AI agents}, over many rounds and across more than a hundred sub-agent assignments.
\textbf{The human author's role was fourfold}: to set the target, to hold the honesty line, to
refuse to call an unproved statement a theorem, and to demand, in every round, an answer to the
question \emph{who is opposing this}. The last of these is the one that mattered most, because it
converted each round from a claim into a contest.
\end{claimbox}

Concretely:

\begin{itemize}
\item \textbf{All thirty-eight ledger entries were caught by independent review} --- AI auditing
AI, with the reviewer separated from the author of the claim. No entry was caught by the agent that
made the error.
\item \textbf{Two of them were caught during the writing of this manuscript}: corpus items 35 and 36
(rows 26 and 27 of Table~\ref{tab:ledger}) --- a negative $r_0$ and an over-broad
asymptotic-perturbation estimate. Both were introduced in earlier
rounds and survived until this section was drafted.
\item \textbf{The downgrade of \code{SG1-1} from \LT{} to \LC{}} was reported by its own author,
after a computability reading broke one step of the argument.
\end{itemize}

\begin{warnbox}
\textbf{Framing, stated explicitly.} The sentence that carries this subsection is not ``AI agents
produced correct mathematics'' and must not be read as a claim of credit for the result. It is:
\emph{the reliability of this paper does not come from AI agents being unable to err; it comes from
the requirement that every error be caught by an independent path.} The thirty-eight ledger entries
are evidence for the second sentence and evidence against the first: an error process that
produced thirty-eight caught collisions is not an error process that can be trusted to have
produced no uncaught ones. The human author's contribution is therefore best described as the
maintenance of that requirement, and the paper's claim to reliability is a claim about the
requirement, not about the agents.
\end{warnbox}

\noindent We also record the negative side of the disclosure: no step of the reduction, no window
correction, and no grade assignment in this paper was accepted on the authority of a single agent,
and the places where a single agent's conclusion was allowed to stand are exactly the places marked
\LE{} --- measured, not proved.

\appendix

\section{Reproducibility}
\label{sec:appendix-repro}

\noindent\emph{This appendix is the reproducibility section of the merged draft; it prints as
Appendix~A because the source issues \code{\textbackslash{}appendix} here.}

Everything numerical, formal, and enumerative in this paper is regenerable from the artefacts
listed here. Paths of all scripts and logs are given in their respective \code{*.md} files, which
are the authoritative index; the summary below is for orientation only.

\subsection*{Numerical measurement}

\begin{itemize}
\item \textbf{Algorithm}: \code{njit} heat bath, single long chain, thermalized once, no
re-thermalization between volumes.
\item \textbf{Parameters}: $\beta = 2.5$, $L_s = 6$ and $L_s = 8$, $L_t = 12$, $12000$
configurations on the shared chain.
\item \textbf{Cost}: $2.223$~ms/sweep at $L = 6$ and $5.390$~ms/sweep at $L = 8$, measured on the
reference machine; the long-chain design means the $L = 8$ run is the binding cost.
\item \textbf{Outputs}: $w_1 = 0.3427 \pm 0.0119$ and $0.0168 \pm 0.0142$ for the two operators
(\S\ref{sec:measurement}); the sampler bug that invalidates four earlier sign determinations is documented with its
fix in the same directory.
\end{itemize}

\subsection*{Formal verification}

\begin{itemize}
\item \textbf{File}: \code{M02-CBAR.lean}.
\item \textbf{Content}: $25$ declarations, \code{sorry = 0}.
\item \textbf{Axioms}: \code{[propext, Classical.choice, Quot.sound]} --- no axiom beyond the three
permitted by the grade rules stated on the title page.
\item \textbf{Scope}: the isolation statement and its supporting lemmas; the reduction itself is
not formalised.
\end{itemize}

\subsection*{Exhaustive enumeration}

\begin{itemize}
\item $a_7 = 189\,215\,964$, exact, by complete enumeration rather than by extrapolation.
\item $b(\Gamma) \le 14n + 6$, and the constant is sharp.
\item $t_n \le B^n$ with $B = 19^{19}/18^{18} = 50.2820906$.
\end{itemize}

\subsection*{Provenance}

Every number quoted above appears with its source file at the point of use in
\S\ref{sec:established}--\S\ref{sec:ledger}. Where a
number is a judgement rather than a measurement --- the $0\%$ of \S\ref{sec:verdict-clay}, the $20$
in ``twenty
routes'', the count of thirty-eight ledger entries --- it is labelled as such at the point of use.
The provenance rule stated on the title page is the one rule this paper does not relax in its own
favour.


% =====================================================================
%  References (thebibliography so that \cite works)
% =====================================================================
\clearpage
\phantomsection

Records were checked against arXiv, Crossref and the publishers' tables of contents; where a record
could not be confirmed from two independent sources the uncertainty is stated in the entry rather
than smoothed over. Three corrections were made in this third edition: the DOI of
\cite{mrs} pointed at an unrelated fluid-dynamics paper and has been replaced; the year of
\cite{gopfert} was 1982, not 1981; and the volume-and-page string attached to \cite{balaban} mixed a
separate paper into the \emph{Large field} series, which is now listed on its own. Where an entry
carries a DOI or an arXiv identifier, that identifier is the checkable part of the record.

{\small
\setlength{\itemsep}{1.2pt}
\begin{thebibliography}{99}

\bibitem{jaffe} A.~Jaffe, E.~Witten, \emph{Quantum Yang--Mills Theory}, Clay Mathematics
  Institute.

\bibitem{sw} R.~F.~Streater, A.~S.~Wightman, \emph{PCT, Spin and Statistics, and All That},
  W.~A.~Benjamin, New York, 1964; reprinted in \emph{Princeton Landmarks in Physics},
  Princeton University Press, 2000.

\bibitem{os1973} Osterwalder--Schrader, Comm.\ Math.\ Phys.\ \textbf{31} (1973) 83.

\bibitem{osterwalder-seiler} K.~Osterwalder, E.~Seiler, \emph{Gauge field theories on a
  lattice}, Ann.\ Phys.\ \textbf{110} (1978) 440--471, DOI 10.1016/0003-4916(78)90039-8.

\bibitem{szz} H.~Shen, R.~Zhu, X.~Zhu, \emph{A stochastic analysis approach to lattice
  Yang--Mills at strong coupling}, Comm.\ Math.\ Phys.\ \textbf{400} (2023) 805--851,
  DOI 10.1007/s00220-022-04609-1, arXiv:2204.12737 [math.PR].

\bibitem{kp} R.~Koteck\'y, D.~Preiss, \emph{Cluster expansion for abstract polymer models},
  Comm.\ Math.\ Phys.\ \textbf{103} (1986) 491--498, DOI 10.1007/BF01211762.

\bibitem{seiler} Seiler, Lecture Notes in Physics \textbf{159} (1982).

\bibitem{bgdf} D.~C.~Brydges, C.~Giffen, B.~Durhuus, J.~Fr\"ohlich, \emph{Surface
  representations of Wilson loop expectations in lattice gauge theory}, Nucl.\ Phys.\ B
  \textbf{275} (1986) 459--487, DOI 10.1016/0550-3213(86)90609-7.

\bibitem{gjs} Glimm--Jaffe--Spencer, Ann.\ Math.\ \textbf{100} (1974) 585--632.

\bibitem{fo} J.~Feldman, K.~Osterwalder, \emph{The Wightman axioms and the mass gap for
  weakly coupled $(\phi^4)_3$ quantum field theories}, Ann.\ Physics \textbf{97} (1976)
  80--135.

\bibitem{balaban} T.~Ba{\l}aban, \emph{Renormalization group approach to lattice gauge field
  theories.\ I}, Comm.\ Math.\ Phys.\ \textbf{109} (1987) 249--301, DOI 10.1007/BF01215223;
  \emph{II.\ Cluster expansions}, Comm.\ Math.\ Phys.\ \textbf{116} (1988) 1--22,
  DOI 10.1007/BF01239022; \emph{Large field renormalization.\ I.\ The basic step of the
  R operation}, Comm.\ Math.\ Phys.\ \textbf{122} (1989) 175--202; \emph{II.\ Localization,
  exponentiation, and bounds for the R operation}, Comm.\ Math.\ Phys.\ \textbf{122} (1989)
  355--392.

\bibitem{balaban-conv} T.~Ba{\l}aban, \emph{Convergent renormalization expansions for
  lattice gauge theories}, Comm.\ Math.\ Phys.\ \textbf{119} (1988) 243--285,
  DOI 10.1007/BF01217741.

\bibitem{mrs} J.~Magnen, V.~Rivasseau, R.~S\'en\'eor, \emph{Construction of $\mathrm{YM}_4$
  with an infrared cutoff}, Comm.\ Math.\ Phys.\ \textbf{155} (1993) 325--383,
  DOI 10.1007/BF02097397.

\bibitem{kl} K\"all\'en, \emph{Helv.\ Phys.\ Acta} \textbf{25} (1952) 417; Lehmann,
  \emph{Nuovo Cim.} \textbf{11} (1954) 342.

\bibitem{prekopa} Pr\'ekopa, Acta Sci.\ Math.\ Szeged \textbf{32} (1971) 301--316.

\bibitem{bl} Brascamp--Lieb, J.\ Funct.\ Anal.\ \textbf{22} (1976) 366.

\bibitem{savvidy} Savvidy, Phys.\ Lett.\ B \textbf{71} (1977) 133.

\bibitem{no} Nielsen--Olesen, Nucl.\ Phys.\ B \textbf{144} (1978) 376.

\bibitem{leutwyler} Leutwyler, Nucl.\ Phys.\ B \textbf{179} (1981) 129.

\bibitem{luscher77} M.~L\"uscher, \emph{Construction of a selfadjoint, strictly positive
  transfer matrix for Euclidean lattice gauge theories}, Comm.\ Math.\ Phys.\ \textbf{54}
  (1977) 283--292, DOI 10.1007/BF01614090; \emph{Some analytic results concerning the mass
  spectrum of Yang--Mills gauge theories on a torus}, Nucl.\ Phys.\ B \textbf{219} (1983)
  233--261.

\bibitem{singer} Singer, Comm.\ Math.\ Phys.\ \textbf{60} (1978) 7.

\bibitem{kato} Kato, \emph{Perturbation Theory for Linear Operators} (1976), VII \S3.

\bibitem{reed-simon} Reed--Simon IV, \S XIII.1.

\bibitem{vnw} von Neumann--Wigner, Phys.\ Z.\ \textbf{30} (1929) 465--467.

\bibitem{morningstar} Morningstar--Peardon, hep-lat/9901004.

\bibitem{bbs} Bauerschmidt--Brydges--Slade, Lecture Notes in Mathematics \textbf{2242} (2019).

\bibitem{sinclair} Sinclair, Combin.\ Probab.\ Comput.\ \textbf{1} (1992) 351.

\bibitem{js} Jerrum--Sinclair, SIAM J.\ Comput.\ \textbf{18} (1989) 1149.

\bibitem{feller} Feller, Ann.\ of Math.\ \textbf{55} (1952) 468.

\bibitem{griffiths} Griffiths, J.\ Math.\ Phys.\ \textbf{8} (1967) 478.

\bibitem{ksherman} Kelly--Sherman, J.\ Math.\ Phys.\ \textbf{9} (1968) 466.

\bibitem{fkg} Fortuin--Kasteleyn--Ginibre, Comm.\ Math.\ Phys.\ \textbf{22} (1971) 89--103.

\bibitem{ginibre} Ginibre, Comm.\ Math.\ Phys.\ \textbf{16} (1970) 310.

\bibitem{es} Efron--Stein, Ann.\ Statist.\ \textbf{9} (1981) 586.

\bibitem{royen} Royen, Far East J.\ Theor.\ Stat.\ \textbf{48} (2014) 139.

\bibitem{vrs} von Renesse--Sturm, Comm.\ Pure Appl.\ Math.\ \textbf{58} (2005) 923.

\bibitem{gsv} Guo--Shamai--Verd\'u, IEEE Trans.\ Inform.\ Theory \textbf{51} (2005) 1261.

\bibitem{bbl} Borcea--Br\"and\'en--Liggett, J.\ Amer.\ Math.\ Soc.\ \textbf{22} (2009)
  521--567.

\bibitem{cornwall} Cornwall, Phys.\ Rev.\ D \textbf{26} (1982) 1453.

\bibitem{temple} Temple, Proc.\ Roy.\ Soc.\ A \textbf{119} (1928) 276--293.

\bibitem{gopfert} M.~G\"opfert, G.~Mack, \emph{Proof of confinement of static quarks in
  3-dimensional $U(1)$ lattice gauge theory for all values of the coupling constant},
  Comm.\ Math.\ Phys.\ \textbf{82} (1982) 545--606, DOI 10.1007/BF01961240.

\bibitem{polchinski} J.~Polchinski, \emph{Scale and conformal invariance in quantum field
  theory}, Nucl.\ Phys.\ B \textbf{303} (1988) 226--236.

\bibitem{dksst} A.~Dymarsky, Z.~Komargodski, A.~Schwimmer, S.~Theisen, \emph{On scale and
  conformal invariance in four dimensions}, JHEP \textbf{10} (2015) 171,
  DOI 10.1007/JHEP10(2015)171, arXiv:1309.2921 [hep-th].

\bibitem{komargodski-schwimmer} Z.~Komargodski, A.~Schwimmer, \emph{On renormalization group
  flows in four dimensions}, JHEP \textbf{12} (2011) 099, arXiv:1107.3987 [hep-th].

\bibitem{banks-zaks} T.~Banks, A.~Zaks, \emph{On the phase structure of vector-like gauge
  theories with massless fermions}, Nucl.\ Phys.\ B \textbf{196} (1982) 189--204,
  DOI 10.1016/0550-3213(82)90035-9.

\bibitem{caswell} W.~E.~Caswell, \emph{Asymptotic behavior of non-Abelian gauge theories to
  two-loop order}, Phys.\ Rev.\ Lett.\ \textbf{33} (1974) 244.

\bibitem{vrvl} T.~van Ritbergen, J.~A.~M.~Vermaseren, S.~A.~Larin, \emph{The four-loop beta
  function in quantum chromodynamics}, Phys.\ Lett.\ B \textbf{400} (1997) 379--384,
  hep-ph/9701390.

\bibitem{hofman-maldacena} D.~M.~Hofman, J.~Maldacena, \emph{Conformal collider physics:
  energy and charge correlations}, JHEP \textbf{0805} (2008) 012, arXiv:0803.1467 [hep-th].

\bibitem{hofman-li} D.~M.~Hofman, D.~Li, E.~Meltzer, D.~Poland, F.~Rej\'on-Barrera, \emph{A
  proof of the conformal collider bounds}, JHEP \textbf{06} (2016) 111,
  DOI 10.1007/JHEP06(2016)111, arXiv:1603.03771 [hep-th].

\bibitem{vsmekal} L.~von Smekal, A.~Hauck, R.~Alkofer, \emph{Infrared behavior of gluon and
  ghost propagators in Landau gauge QCD}, Phys.\ Rev.\ Lett.\ \textbf{79} (1997) 3591--3594,
  DOI 10.1103/PhysRevLett.79.3591.

\bibitem{fischer-alkofer} C.~S.~Fischer, R.~Alkofer, \emph{Nonperturbative propagators, running
  coupling and dynamical mass generation in Landau-gauge QCD}, Phys.\ Lett.\ B \textbf{536}
  (2002) 177--184, hep-ph/0202202.

\bibitem{huber} M.~Q.~Huber, \emph{Nonperturbative properties of Yang--Mills theories},
  Phys.\ Rept.\ \textbf{879} (2020) 1--92, arXiv:1808.05227 [hep-ph].

\bibitem{gribov} V.~N.~Gribov, \emph{Quantization of non-Abelian gauge theories}, Nucl.\ Phys.\
  B \textbf{139} (1978) 1--19.

\bibitem{zwanziger} D.~Zwanziger, \emph{Local and global gauge fixing in non-Abelian gauge
  theories}, Nucl.\ Phys.\ B \textbf{323} (1989) 513--544.

\bibitem{aoki} Y.~Aoki et al., \emph{Light composite scalar in twelve-flavor QCD on the
  lattice}, Phys.\ Rev.\ Lett.\ \textbf{111} (2013) 162001, arXiv:1305.6006 [hep-lat].

\bibitem{luscher-flow} M.~L\"uscher, \emph{Properties and uses of the Wilson flow in lattice
  QCD}, JHEP \textbf{1008} (2010) 071, arXiv:1006.4518 [hep-lat].

\bibitem{makeenko-migdal} Yu.~M.~Makeenko, A.~A.~Migdal, \emph{Exact equation for the loop
  average in multicolor QCD}, Phys.\ Lett.\ B \textbf{88} (1979) 135.

\bibitem{eguchi-kawai} T.~Eguchi, H.~Kawai, \emph{How to calculate the Wilson loop in the
  large-$N$ limit}, Phys.\ Rev.\ Lett.\ \textbf{48} (1982) 1063.

\bibitem{polchinski-strominger} J.~Polchinski, A.~Strominger, \emph{Effective string theory},
  Phys.\ Rev.\ Lett.\ \textbf{67} (1991) 1681--1684.

\bibitem{luscher-weisz} M.~L\"uscher, P.~Weisz, \emph{Quark confinement and the bosonic string},
  JHEP \textbf{0207} (2002) 049, hep-lat/0207003.

\bibitem{zamolodchikov} A.~B.~Zamolodchikov, \emph{Irreversibility of the flux of the
  renormalization group in a 2D field theory}, JETP Lett.\ \textbf{43} (1986) 730.

\bibitem{osborn} H.~Osborn, \emph{Implications of conformal invariance for quantum field
  theories}, Nucl.\ Phys.\ B \textbf{363} (1991) 486.

\bibitem{maldacena} J.~M.~Maldacena, \emph{The large $N$ limit of superconformal field theories
  and supergravity}, Adv.\ Theor.\ Math.\ Phys.\ \textbf{2} (1998) 231.

\bibitem{gkp} S.~S.~Gubser, I.~R.~Klebanov, A.~M.~Polyakov, \emph{Gauge theory correlators from
  non-critical string theory}, Phys.\ Lett.\ B \textbf{428} (1998) 105.

\bibitem{witten} E.~Witten, \emph{Anti de Sitter space and holography},
  Adv.\ Theor.\ Math.\ Phys.\ \textbf{2} (1998) 253.

\bibitem{chatterjee} S.~Chatterjee, \emph{Yang--Mills for probabilists}, arXiv:1801.03247
  [math.PR] (lecture notes).

\bibitem{wilson} K.~G.~Wilson, \emph{Confinement of quarks}, Phys.\ Rev.\ D \textbf{10} (1974)
  2445--2459.

\bibitem{creutz} M.~Creutz, \emph{Monte Carlo study of quantized SU(2) gauge theory}, Phys.\
  Rev.\ D \textbf{21} (1980) 2308.

\bibitem{luscher-stepscaling} M.~L\"uscher, R.~Sommer, P.~Weisz, U.~Wolff, \emph{A new method to
  compute the running coupling in asymptotically free theories}, Nucl.\ Phys.\ B \textbf{413}
  (1994) 481--502, hep-lat/9309005.

\bibitem{witten-thermal} E.~Witten, \emph{Anti-de Sitter space, thermal phase transition, and
  confinement in gauge theories}, Adv.\ Theor.\ Math.\ Phys.\ \textbf{2} (1998) 505--532,
  hep-th/9803131.

\bibitem{polchinski-strassler} J.~Polchinski, M.~J.~Strassler, \emph{Hard scattering and
  gauge/string duality}, Phys.\ Rev.\ Lett.\ \textbf{88} (2002) 031601, hep-th/0109174.

\bibitem{sakai-sugimoto} T.~Sakai, S.~Sugimoto, \emph{Low energy hadron physics in holographic
  QCD}, Prog.\ Theor.\ Phys.\ \textbf{113} (2005) 843--882, hep-th/0412141.

\bibitem{elshowk} S.~El-Showk, M.~F.~Paulos, D.~Poland, S.~Rychkov, D.~Simmons-Duffin,
  A.~Vichi, \emph{Solving the 3d Ising model with the conformal bootstrap}, Phys.\ Rev.\ D
  \textbf{86} (2012) 025022.

\bibitem{prv} D.~Poland, S.~Rychkov, A.~Vichi, \emph{The conformal bootstrap: theory, numerical
  techniques, and applications}, Rev.\ Mod.\ Phys.\ \textbf{91} (2019) 015002.

\bibitem{abs} I.~Aniceto, G.~Ba\c{s}ar, R.~Schiappa, \emph{A primer on resurgent transseries and
  their asymptotics}, Phys.\ Rept.\ \textbf{809} (2019) 1--135.

\bibitem{thooft-planar} G.~'t~Hooft, \emph{A planar diagram theory for strong interactions},
  Nucl.\ Phys.\ B \textbf{72} (1974) 461.

\end{thebibliography}
}


\end{document}
